\documentclass[11pt, oneside]{article}

\usepackage{amsmath, amsthm, amssymb, amsfonts}
\usepackage{mathtools}
\usepackage{geometry}
\usepackage{hyperref}
\usepackage{booktabs}
\usepackage{array}
\usepackage{xcolor}
\usepackage{titlesec}
\usepackage{abstract}
\usepackage{microtype}
\usepackage{setspace}
\usepackage[numbers, sort&compress]{natbib}

\geometry{
  paper=a4paper,
  top=2.5cm, bottom=2.5cm,
  left=3cm, right=3cm
}

\onehalfspacing

\hypersetup{
  colorlinks=true,
  linkcolor=black,
  citecolor=black,
  urlcolor=black
}

% Theorem environments
\newtheorem{theorem}{Theorem}[section]
\newtheorem{lemma}[theorem]{Lemma}
\newtheorem{proposition}[theorem]{Proposition}
\newtheorem{corollary}[theorem]{Corollary}
\newtheorem{definition}[theorem]{Definition}
\newtheorem{assumption}[theorem]{Assumption}
\newtheorem{remark}[theorem]{Remark}

% Preliminary section styling
\newcommand{\preliminary}[1]{%
  \colorbox{yellow!20}{\parbox{\dimexpr\linewidth-2\fboxsep}{%
    \textbf{\textcolor{red!70!black}{[PRELIMINARY OUTLINE --- NOT YET DRAFTED]}}\\[4pt]
    #1
  }}%
}

% Math operators
\DeclareMathOperator{\supp}{supp}
\DeclareMathOperator{\vol}{Vol}
\DeclareMathOperator{\interior}{int}
\DeclareMathOperator{\diam}{diam}
\newcommand{\norm}[1]{\left\|#1\right\|}
\newcommand{\abs}[1]{\left|#1\right|}
\newcommand{\inner}[2]{\langle #1, #2 \rangle}
\newcommand{\Xmod}{\mathcal{X}/{\sim}}
\newcommand{\Xquot}{\mathcal{X}/G}
\newcommand{\Ftil}{\widetilde{\Pi}}

\title{%
  \Large\textbf{Yarncrawler: World-State Reconstruction}\\[6pt]
  \large\textbf{as Constraint Closure over a Generative Field}
}

\author{Flyxion}
\date{\today}

\begin{document}

\maketitle

\begin{abstract}
This paper introduces Yarncrawler as a framework for reconstructing physical world-states
from multi-modal sensor observations under continuous dynamical constraint. The central
object is a consistency operator $\mathcal{C}$ acting on the quotient configuration space
$\mathcal{X}/{\sim}$ of coupled Relativistic Scalar-Vector Plenum field triples
$(\Phi, \mathbf{v}, S)$. Each sensor modality defines a projection operator
$\Pi_i : \mathcal{X} \to \mathcal{Y}_i$; each observation $y_i$ constrains the set of
field configurations consistent with what the sensor reports. The consistency operator
searches for the unique admissible configuration whose forward projections are compatible
with all observations simultaneously.

The main result---the Yarncrawler Identifiability Theorem---establishes that $\mathcal{C}$
has a unique fixed point if and only if the induced map
$\widetilde{\Pi} : \Xmod \to \prod_i \mathcal{Y}_i$ is injective on the feasible set
$\mathcal{F} = \Omega_{\mathrm{obs}} \cap \Gamma$. The Dynamical Disambiguation Theorem
extends this to realistic sensing regimes where static observations underdetermine the
world-state: trajectory fragments constrain the tangent bundle of the admissible manifold,
removing degrees of freedom that no static projection can touch.

The paper further establishes that the variational consistency framework and the
sheaf-theoretic repair framework---in which sensor modalities define local charts, observations
define local sections, and closure corresponds to the existence of a unique global
section---are formally equivalent. There exists a functor
$\Psi : \mathbf{Proj}(\mathcal{X}) \to \mathbf{Sh}(\Xmod)$ under which the consistency
operator is the variational realization of the sheaf repair operator, projection degeneracy
corresponds to non-trivial $H^1$, and entropy collapse corresponds to cohomology
trivialization. The RSVP entropy field $S$ is identified as the local density of
cohomological obstruction; its minimization under $\mathcal{C}$ is the physical expression
of sheaf repair.

Deployment at planetary, regional, and site scales is unified by treating each scale as a
different geometry of the feasible set under varying noise profiles. The framework is
implemented as a sensor fusion architecture combining learned trajectory operators with
physics-constrained forward models, applied to coupled ecosystem and civilization dynamics
calibrated from before-and-after observational pairs across modalities including infrared,
LiDAR, RGB, audio, and photonic strip sensors.
\end{abstract}

\tableofcontents
\newpage

%% ============================================================
\section{Introduction}
%% ============================================================

\subsection{The Problem}

Given a family of projection operators $\{\Pi_i\}_{i \in \mathcal{I}}$ mapping an unknown
world-state $[X] \in \Xmod$ to observable quantities $y_i \in \mathcal{Y}_i$, and given a
dynamical law governing how world-states evolve, the central question of this paper is:

\begin{quote}
\textit{Under what conditions does there exist a unique trajectory in the quotient
configuration space $\Xmod$ consistent with all projections simultaneously and admissible
under the dynamics?}
\end{quote}

This is not a modeling question. It is an identifiability question. The paper does not
propose a method for approximating world-states from sensor data---it establishes the
conditions under which a unique world-state is \emph{demanded} by the data, and proves
that when those conditions hold, a specific operator converges to it.

The setting is the following. The world-state is a coupled field triple
$X = (\Phi, \mathbf{v}, S)$---scalar density, vector flow, and entropy
production---evolving over a compact spatial domain under the Relativistic Scalar-Vector
Plenum (RSVP) field equations. The projections $\Pi_i$ correspond to sensor modalities:
infrared, LiDAR, RGB, audio, photonic strips. Each observation $y_i$ constrains which
field configurations are consistent with what the sensor reports. The consistency operator
$\mathcal{C}$ searches for the configuration that minimizes total projection residual
subject to physical admissibility. The dynamical admissibility set $\Gamma$ restricts
attention to configurations that lie on trajectories the world could have taken. The
feasible set $\mathcal{F} = \Omega_{\mathrm{obs}} \cap \Gamma$ is the intersection of all
constraints. Closure---the condition that $|\mathcal{F}| = 1$---is the condition that the
world-state is uniquely determined.

\subsection{The Central Result}

The paper proves the following.

The problem of multi-modal consistency---finding a world-state whose forward projections
are compatible with all available observations---is equivalent to the existence of a unique
global section of a sheaf $\mathcal{S}_\Pi$ induced by the projection family over a cover
of $\Xmod$.

More precisely: the consistency operator $\mathcal{C}$ is the variational realization of
the sheaf repair operator. Its fixed points are global sections. Its failure modes---
projection degeneracy and regularizer flatness---are cohomological obstructions $H^1 \ne 0$
and degenerate $H^0$. Adding sensors refines the cover and kills $H^1$ classes.
Strengthening the dynamical constraint extends the cover into the temporal direction and
kills residual ambiguity that spatial sensing cannot reach. Closure is the condition
$H^0 = 1$, $H^1 = 0$: the sheaf has exactly one global section.

This is the Yarncrawler Identifiability Theorem and its dynamical extension, proved in
Sections \ref{sec:identifiability} and \ref{sec:disambiguation} and their appendices.

\subsection{The Duality}

The framework has two faces, both present in this paper.

The \emph{forward} direction takes a world-state as given and asks how a system maintains
coherence under perturbation: how it detects tears in its semantic fabric, applies repair
morphisms that trivialize \v{C}ech cocycles, exports entropy through stigmergic
reinforcement, and sustains homeorhetic flows rather than static equilibria.

The \emph{inverse} direction takes sensor observations as given and asks what world-state
they structurally demand. Given what the sensors report, what field configuration is forced?

These are not complementary perspectives on the same phenomenon. They are dual realizations
of the same mathematical object. Both directions study the sheaf $\mathcal{S}_\Pi$ over the
quotient configuration space $\Xmod$. The forward direction studies the internal dynamics
of sections within $\mathcal{S}_\Pi$---how they evolve, heal, and accumulate into global
coherence. The inverse direction studies the conditions under which a unique global section
exists and the operator that converges to it. The RSVP field triple $(\Phi, \mathbf{v}, S)$
parametrizes the stalks in both directions. The Yarncrawler framework is the unified theory
of that sheaf and its dynamics.

The object common to both is neither the operator nor the sheaf taken separately. It is
the equivalence between them---the functor
$\Psi : \mathbf{Proj}(\mathcal{X}) \to \mathbf{Sh}(\Xmod)$ proved in
Section~\ref{sec:equivalence}, which shows that a projection family and a sheaf over a
configuration space are the same structure viewed from outside and inside respectively.

\subsection{The Consequence}

When closure is achieved---when $|\mathcal{F}| = 1$---three things happen simultaneously
that were previously three separate conditions.

The projection residuals collapse to within noise tolerance:
$\|\Pi_i(X^*) - y_i\| \le \epsilon_i$ for all $i$. The cohomology of the induced sheaf
trivializes: $H^1(\mathcal{U}, \mathcal{S}_\Pi) = 0$ and $\dim H^0 = 1$. The RSVP entropy
field $S$ is minimized subject to admissibility: no further constraint can be satisfied
that is not already satisfied.

These are not three criteria that must be checked independently. Under the
Sheaf-Variational Equivalence of Section~\ref{sec:equivalence}, they are one criterion
expressed in three languages---variational, cohomological, and field-theoretic. The entropy
of the RSVP field at the fixed point is the local density of cohomological obstruction,
and its minimization is the physical expression of sheaf repair.

The closed state $[X^*]$ is therefore not a best estimate. It is the unique trajectory
consistent with all projections and admissible dynamics---the world-state that no
alternative history can explain away.

\subsection{Structure of the Paper}

Section~\ref{sec:equivalence} establishes the Sheaf-Variational Equivalence. Sections
\ref{sec:worldstate} and \ref{sec:projections} develop the prerequisite machinery.
Section~\ref{sec:feasible} defines the feasible set geometrically.
Sections~\ref{sec:identifiability} and \ref{sec:disambiguation} contain the core theorems.
Section~\ref{sec:inversion} develops the inversion layer architecture.
Section~\ref{sec:closure} defines closure operationally and proves equivalence of the three
closure criteria. Section~\ref{sec:deployment} unifies planetary, regional, and site-scale
deployment. Section~\ref{sec:related} positions the framework against prior work.
Section~\ref{sec:epistemology} draws the epistemological consequence. Appendices
\ref{app:frechet}--\ref{app:operators} provide the functional-analytic foundations.

%% ============================================================
\section{The Sheaf-Variational Equivalence}
\label{sec:equivalence}
%% ============================================================

\subsection{Two Descriptions of the Same Object}

The Yarncrawler framework has two faces. The first describes a self-refactoring system
that maintains coherence by repairing tears in its semantic fabric: local sections are
parsed, inconsistencies are detected as \v{C}ech cocycles, and repair morphisms trivialize
them until a global section exists. The second describes a consistency operator that
searches for a field configuration whose forward projections are compatible with all
available sensor observations simultaneously. These descriptions appear to operate at
different levels of abstraction---one categorical and topological, the other variational
and analytic. This section proves they are the same description.

The identification is not a metaphor. There is a precise functor between the category of
sheaves of semantic sections over a cover and the category of projection families on the
quotient configuration space $\Xmod$. Under this functor, every concept in the
sheaf-theoretic account has an exact counterpart in the variational account, and vice
versa.

\subsection{The Sheaf-Theoretic Account}

Let $\mathcal{M}$ be a topological space equipped with an open cover
$\mathcal{U} = \{U_i\}_{i \in \mathcal{I}}$. A presheaf $\mathcal{S}$ over $\mathcal{U}$
assigns to each $U_i$ a category $\mathcal{S}(U_i)$ of local sections, and to each
inclusion $U_j \subseteq U_i$ a restriction functor
$\rho_{ij} : \mathcal{S}(U_i) \to \mathcal{S}(U_j)$.

\begin{definition}[Gluing Condition]
A family of local sections $\{s_i \in \mathcal{S}(U_i)\}_{i \in \mathcal{I}}$ satisfies
the \emph{gluing condition} if for all $i, j \in \mathcal{I}$:
\[
\rho_{U_i, U_i \cap U_j}(s_i) = \rho_{U_j, U_i \cap U_j}(s_j) \quad \text{on }
U_i \cap U_j.
\]
A presheaf in which every compatible family of local sections admits a unique global
section $s \in \mathcal{S}(\mathcal{M})$ with $s|_{U_i} = s_i$ is called a \emph{sheaf}.
\end{definition}

\begin{definition}[Semantic Tear and Repair]
A \emph{semantic tear} at overlap $U_i \cap U_j$ is a failure of the gluing condition.
The discrepancy defines a \v{C}ech 1-cochain
$c_{ij} = \rho_{ij}(s_j) - \rho_{ji}(s_i)$. A \emph{repair} is a morphism in
$\mathcal{S}(U_i \cap U_j)$ that trivializes $c_{ij}$. The space of unresolved tears is
measured by $H^1(\mathcal{U}, \mathcal{S})$.
\end{definition}

\begin{definition}[Sheaf Repair Operator]
The \emph{sheaf repair operator} $\mathcal{R}_{\mathrm{sheaf}}$ maps a family of
inconsistent local sections to a minimally modified family satisfying the gluing condition,
where minimality is measured by a seam cost functional $\mathcal{L}_{\mathrm{seam}}$:
\[
\mathcal{R}_{\mathrm{sheaf}}(\{s_i\}) = \arg\min_{\{s'_i\}} \sum_{i < j}
\mathcal{L}_{\mathrm{seam}}(s'_i|_{U_{ij}}, s'_j|_{U_{ij}}) +
\sum_i \mathcal{L}_{\mathrm{local}}(s'_i, s_i).
\]
Global section existence corresponds to $\mathcal{L}_{\mathrm{seam}} = 0$ at the optimum.
\end{definition}

\subsection{The Variational Account}

The variational account centers on the consistency operator:
\[
\mathcal{C}([X]) = \arg\min_{[X'] \in \mathcal{A}}
\left[ \sum_{i \in \mathcal{I}} \|\Pi_i(X') - y_i\|_{\mathcal{Y}_i}^2
+ \mathcal{R}(X') \right]
\]
where $\Pi_i : \mathcal{X} \to \mathcal{Y}_i$ are projection operators, $y_i$ are
observations, and $\mathcal{R}$ is the RSVP-consistent regularizer encoding admissibility.
A fixed point $[X^*]$ satisfies $\Pi_i(X^*) \approx y_i$ for all $i$ within tolerance
$\epsilon_i$.

\subsection{The Dictionary}

The following correspondence identifies every object in the sheaf-theoretic account with
its exact counterpart in the variational account.

\begin{center}
\begin{tabular}{ll}
\toprule
\textbf{Sheaf-theoretic concept} & \textbf{Variational counterpart} \\
\midrule
Open set $U_i$ in cover $\mathcal{U}$ & Sensor modality $i$ with space $\mathcal{Y}_i$ \\
Local section $s_i \in \mathcal{S}(U_i)$ & Observation $y_i \in \mathcal{Y}_i$ \\
Restriction map $\rho_{ij}$ & Projection operator $\Pi_i$ \\
Gluing condition & Projection consistency $\Pi_i(X) = y_i$ \\
\v{C}ech 1-cochain $c_{ij}$ & Projection residual $\|\Pi_i(X) - y_i\|$ \\
Coboundary (trivialized cocycle) & Residual driven to zero by $\mathcal{C}$ \\
Global section $s \in \mathcal{S}(\mathcal{M})$ & Fixed point $[X^*] \in \Xmod$ \\
Cohomological obstruction $H^1 \ne 0$ & Projection degeneracy \\
Flat direction in $\mathcal{L}_{\mathrm{local}}$ & Regularizer flatness \\
Repair morphism trivializing $c_{ij}$ & One step of the refinement loop \\
Sheaf repair operator $\mathcal{R}_{\mathrm{sheaf}}$ & Consistency operator $\mathcal{C}$ \\
Stigmergic threshold & Contractivity condition on $\mathcal{C}$ \\
Self-maintaining growth & Convergence of refinement loop to $[X^*]$ \\
\bottomrule
\end{tabular}
\end{center}

\subsection{The Sheaf-Variational Equivalence}

\begin{definition}[Category of Projection Families]
Let $\mathbf{Proj}(\mathcal{X})$ be the category whose objects are pairs
$(\mathcal{I}, \{\Pi_i\}_{i \in \mathcal{I}})$ and whose morphisms are inclusions of
index sets compatible with the projection structure.
\end{definition}

\begin{definition}[Category of Covering Sheaves]
Let $\mathbf{Sh}(\mathcal{M})$ be the category whose objects are sheaves $\mathcal{S}$
over covers $\mathcal{U}$ of $\mathcal{M}$, and whose morphisms are morphisms of sheaves
over refinements of covers.
\end{definition}

\begin{theorem}[Sheaf-Variational Equivalence]
\label{thm:equivalence}
There exists a functor
\[
\Psi : \mathbf{Proj}(\mathcal{X}) \longrightarrow \mathbf{Sh}(\Xmod)
\]
such that:
\begin{enumerate}
\item[(i)] $\Psi$ maps each projection family $\{\Pi_i\}_{i \in \mathcal{I}}$ to a sheaf
$\mathcal{S}_\Pi$ over a cover of $\Xmod$ whose local sections are observation-consistent
field configurations.
\item[(ii)] The gluing condition of $\mathcal{S}_\Pi$ is equivalent to the projection
consistency condition $\Pi_i(X) = y_i$ for all $i$.
\item[(iii)] The consistency operator $\mathcal{C}$ is the variational realization of the
sheaf repair operator $\mathcal{R}_{\mathrm{sheaf}}$ under $\Psi$.
\item[(iv)] The fixed point $[X^*]$ of $\mathcal{C}$ corresponds under $\Psi$ to the
unique global section of $\mathcal{S}_\Pi$.
\item[(v)] The cohomological obstruction $H^1(\mathcal{U}, \mathcal{S}_\Pi) = 0$ if and
only if the projection family is identifying on $\mathcal{F}$.
\end{enumerate}
\end{theorem}

\begin{proof}
\textbf{Construction of $\Psi$.} Given a projection family
$(\mathcal{I}, \{\Pi_i\})$, construct the cover
$\mathcal{U} = \{U_i\}_{i \in \mathcal{I}}$ of $\Xmod$ where
\[
U_i = \{ [X] \in \mathcal{A} \mid \|\Pi_i(X) - y_i\|_{\mathcal{Y}_i} \le \epsilon_i \}.
\]
Define the sheaf $\mathcal{S}_\Pi$ by assigning to each $U_i$:
\[
\mathcal{S}_\Pi(U_i) = U_i \cap \mathcal{A}.
\]
For inclusions $U_j \subseteq U_i$, the restriction map $\rho_{ij}$ is the inclusion
$U_j \cap \mathcal{A} \hookrightarrow U_i \cap \mathcal{A}$.

\textbf{Part (i).} By construction, $\mathcal{S}_\Pi(U_i)$ consists of admissible
configurations whose $i$-th projection is within tolerance of $y_i$. $\checkmark$

\textbf{Part (ii).} A family $\{[X_i] \in \mathcal{S}_\Pi(U_i)\}$ satisfies the gluing
condition iff $[X_i] = [X_j]$ on $U_i \cap U_j$ for all $i,j$, which is precisely the
condition that a single $[X] \in \mathcal{A}$ satisfies $\Pi_i(X) = y_i$ for all $i$
simultaneously. $\checkmark$

\textbf{Part (iii).} The objective of $\mathcal{R}_{\mathrm{sheaf}}$ translates as:
\[
\sum_{i<j} \mathcal{L}_{\mathrm{seam}} + \mathcal{L}_{\mathrm{local}}
\;\longleftrightarrow\;
\sum_i \|\Pi_i(X) - y_i\|^2 + \mathcal{R}(X),
\]
which is exactly the objective of $\mathcal{C}$. $\checkmark$

\textbf{Part (iv).} A fixed point $[X^*]$ of $\mathcal{C}$ satisfies
$[X^*] \in U_i$ for all $i$, giving a local section $s_i = [X^*]|_{U_i}$ that trivially
satisfies the gluing condition. Uniqueness follows from Theorem~\ref{thm:identifiability}.
$\checkmark$

\textbf{Part (v).} $H^1(\mathcal{U}, \mathcal{S}_\Pi)$ measures failure of compatible
local sections to glue into a global section, which is exactly projection degeneracy by
part (ii). $\checkmark$

\textbf{Functoriality.} Adding sensors corresponds to refining the cover; the refinement
map defines a natural transformation establishing functoriality. $\square$
\end{proof}

\subsection{Consequences of the Equivalence}

\textbf{Consequence 1: The two failure modes are cohomological.} Projection degeneracy is
$H^1 \ne 0$; regularizer flatness is degenerate $H^0$. Adding sensors can kill $H^1$
classes; strengthening the regularizer thins the stalks and kills degenerate $H^0$.

\textbf{Consequence 2: Dynamical disambiguation is cohomology in time.} The dynamical
admissibility set $\Gamma$ extends the cover into the temporal direction. Trajectory
separation kills residual $H^1$ classes that spatial sensors cannot. The full feasible set
$\mathcal{F} = \Omega_{\mathrm{obs}} \cap \Gamma$ is the global section of a sheaf over a
spatio-temporal cover.

\textbf{Consequence 3: Forward and inverse are dual.} The forward direction studies the
repair dynamics of $\mathcal{S}_\Pi$ from inside---how sections evolve under stigmergic
reinforcement and homeorhetic flow. The inverse direction studies the existence and
uniqueness of the global section from outside. Both study the same sheaf.

\subsection{The Entropy Identification}

The RSVP entropy field $S$ appears in two roles throughout this framework: as the
entropy component of the field triple $(\Phi, \mathbf{v}, S)$, and as a measure of
unresolved obstruction. Under the Sheaf-Variational Equivalence these become one.

The observation-consistent set $\Omega_{\mathrm{obs}}$ carries a natural measure induced
by the Fisher information metric of the noise models $\{\mathcal{N}_i\}$ (developed in
Appendix~\ref{app:fisher}). The volume of $\Omega_{\mathrm{obs}}$ under this metric
measures residual uncertainty. Under the equivalence:
\[
\log \vol_{\mathrm{Fisher}}(\Omega_{\mathrm{obs}})
\;\longleftrightarrow\;
\log \dim H^0(\mathcal{U}, \mathcal{S}_\Pi)
\;\longleftrightarrow\;
S_{\mathrm{RSVP}}(X^*).
\]
All three measure residual degrees of freedom after all constraints have been applied.
Closure corresponds to each collapsing simultaneously. The RSVP entropy field is not
merely a bookkeeping device---it is the local density of cohomological obstruction, and
its minimization under $\mathcal{C}$ is the physical expression of sheaf repair.

%% ============================================================
\section{The World-State Manifold and Its Symmetry Group}
\label{sec:worldstate}
%% ============================================================

\subsection{Configuration Space}

The configuration space is
\[
\mathcal{X} = H^1(\Omega) \times H^1(\Omega;\mathbb{R}^3) \times L^2(\Omega)
\]
over a compact domain $\Omega \subset \mathbb{R}^3$ with smooth boundary, equipped with
the product Sobolev norm. The three components have physical interpretations that ground
the abstract framework. The scalar field $\Phi \in H^1(\Omega)$ encodes ecosystem density
and civilizational legitimacy potential: high $\Phi$ at a location corresponds to dense
biomass, concentrated infrastructure, or high resource availability. The vector field
$\mathbf{v} \in H^1(\Omega;\mathbb{R}^3)$ encodes resource flows and infrastructure
transport: its divergence marks sources and sinks, its curl marks circulation patterns.
The entropy field $S \in L^2(\Omega)$ encodes complexity budget and irreversibility: high
$S$ corresponds to high local disorder, biological activity, or unresolved ambiguity in
the system's reconstruction of the world-state.

As a product of Hilbert spaces, $\mathcal{X}$ is itself a separable Hilbert space and
hence a Fréchet manifold. The full Fréchet manifold structure, including the treatment of
extensions to more general field configurations, is developed in
Appendix~\ref{app:frechet}.

\subsection{Symmetry Group}

A symmetry group $G$ acts on $\mathcal{X}$ by three classes of transformations. First,
\emph{gauge redundancies} of the field representation: rephasing and rescaling
transformations that leave all physical observables invariant but change the coordinate
representation of $(\Phi, \mathbf{v}, S)$. Second, \emph{coordinate reparametrizations}
of $\Omega$: smooth diffeomorphisms that relabel spatial positions without altering field
content. Third, \emph{observationally equivalent reconfigurations} of the civilizational
layer: distinct infrastructure configurations or demographic distributions that produce
identical outputs under all available sensor projections.

The action of $G$ is smooth, proper, and isometric with respect to the product Sobolev
metric, so the quotient $\mathcal{X}/G$ is Hausdorff and inherits a natural metric.

Two distinct equivalence relations are in play. $G$-equivalence ($X \sim_G X'$ when
$X' = g \cdot X$ for some $g \in G$) is intrinsic to the field representation.
Observational equivalence ($X \sim X'$ when $\Pi_i(X) = \Pi_i(X')$ for all $i$) is
extrinsic and depends on the projection family. In general $\sim$ is coarser than
$\sim_G$: two configurations may agree on all projections without being related by any
element of $G$. When the projection family is complete in the sense of
Proposition~\ref{prop:completeness}, the two equivalences coincide.

\subsection{The Admissible Set}

Define the admissible set
\[
\mathcal{A} = \{ [X] \in \mathcal{X}/{\sim} \mid \mathcal{R}(X) < \infty \},
\]
where $\mathcal{R}$ is the regularizer defined below. The admissible set consists of
equivalence classes satisfying the RSVP field equations, ecological plausibility, and
civilizational feasibility simultaneously. It is a closed convex subset of the Banach
realization $\mathcal{B} = \mathcal{X}$ (Assumption~\ref{ass:banach}).

Geodesic convexity of $\mathcal{A}$ ensures that admissible interpolations between two
configurations remain admissible. Lower semicontinuity of $\mathcal{R}$ ensures
admissibility is stable under limits. Coercivity ensures that physically implausible
configurations --- those with unbounded field energy or violating conservation laws ---
are excluded rather than merely penalized.

\subsection{The RSVP Field Equations}

The time evolution of $(\Phi, \mathbf{v}, S)$ is governed by the RSVP plenum equations:
\begin{align*}
\partial_t \Phi + \nabla \cdot (\Phi \mathbf{v}) &= 0
\quad \text{(scalar continuity)}, \\
\partial_t \mathbf{v} + (\mathbf{v} \cdot \nabla)\mathbf{v} + \nabla\Phi &= \mathbf{f}
\quad \text{(vector momentum)}, \\
\partial_t S + \nabla \cdot (S\mathbf{v}) &= \sigma(X)
\quad \text{(entropy production)},
\end{align*}
where $\mathbf{f}$ encodes external forcing (ecological or civilizational) and
$\sigma(X) \ge 0$ is the entropy production rate satisfying the second law. Solutions
to these equations form the dynamical admissibility set $\Gamma$ in the zero-noise limit
(Section~\ref{sec:disambiguation}).

\subsection{The Regularizer}

The regularizer encodes coupled closure across all three constraint classes:
\[
\mathcal{R}(X) = \mathcal{R}_{\mathrm{RSVP}}(X)
+ \lambda_{\mathrm{eco}}\,\mathcal{R}_{\mathrm{eco}}(X)
+ \lambda_{\mathrm{civ}}\,\mathcal{R}_{\mathrm{civ}}(X).
\]
$\mathcal{R}_{\mathrm{RSVP}}$ is the squared $L^2$ norm of residuals of the three
field equations above. $\mathcal{R}_{\mathrm{eco}}$ is zero on ecologically admissible
configurations (biomass conservation, succession ordering, watershed feasibility) and
strictly positive off them. $\mathcal{R}_{\mathrm{civ}}$ is zero on civilizationally
admissible configurations (demographic feasibility, infrastructure stress bounds,
resource flow conservation) and strictly positive off them.

The decomposition is notational: $\mathcal{R}$ is a single functional encoding coupled
closure. A configuration satisfying the RSVP equations but violating ecological
admissibility is not merely penalized; it lies outside $\mathcal{A}$ entirely. The
weights $\lambda_{\mathrm{eco}}, \lambda_{\mathrm{civ}} > 0$ are fixed hyperparameters
determining the relative stringency of the three constraint classes.

%% ============================================================
\section{Sensor Modalities as Projection Operators}
\label{sec:projections}
%% ============================================================

\subsection{Purpose of This Section}

Section~\ref{sec:equivalence} established that a projection family $\{\Pi_i\}_{i \in \mathcal{I}}$
and a sheaf $\mathcal{S}_\Pi$ over $\Xmod$ are the same structure viewed from outside and
inside respectively. What remains is to specify what the projection operators actually
are---not as abstract bounded linear maps, but as concrete operators whose domains,
codomains, and kernels are determined by the physics of each sensor modality.

This section defines each projection operator precisely enough that Assumption A1 can be
verified directly. It establishes the RSVP correspondence: the identification of each
modality with a specific component or functional of the field triple $(\Phi, \mathbf{v}, S)$.
And it proves the completeness proposition: the condition under which the projection family
collectively spans the full RSVP triple.

Nothing in this section is motivated by sensor engineering. The operators are defined as
they are because those definitions are what the theorems in Sections~\ref{sec:identifiability}
and \ref{sec:disambiguation} require.

\subsection{The Configuration Space and Its Topology}

Let $\Omega \subset \mathbb{R}^3$ be a compact domain with smooth boundary. The
configuration space is
\[
\mathcal{X} = H^1(\Omega) \times H^1(\Omega; \mathbb{R}^3) \times L^2(\Omega)
\]
with product Sobolev norm $\|X\|_{\mathcal{X}}^2 = \|\Phi\|_{H^1}^2 + \|\mathbf{v}\|_{H^1}^2 + \|S\|_{L^2}^2$.
All projection operators defined below are required to descend to the quotient $\Xmod$:
they must be constant on $G$-orbits.

\subsection{Infrared Projection}

\begin{definition}[Infrared Projection]
The infrared projection $\Pi_{\mathrm{IR}} : \mathcal{X} \to \mathcal{Y}_{\mathrm{IR}}$
is defined by
\[
\Pi_{\mathrm{IR}}(X) = \left( \Phi|_{\partial\Omega},\;
\nabla\Phi \cdot \hat{n}|_{\partial\Omega} \right)
\]
where $\hat{n}$ is the outward unit normal to $\partial\Omega$ and
$\mathcal{Y}_{\mathrm{IR}} = L^2(\partial\Omega) \times L^2(\partial\Omega)$.
\end{definition}

\begin{lemma}
$\Pi_{\mathrm{IR}}$ is a bounded linear operator from $\mathcal{X}$ to $\mathcal{Y}_{\mathrm{IR}}$.
\end{lemma}
\begin{proof}
The trace operator $\Phi \mapsto \Phi|_{\partial\Omega}$ is bounded from $H^1(\Omega)$
to $L^2(\partial\Omega)$ by the Sobolev trace theorem. The normal derivative is bounded
from $H^1(\Omega)$ to $H^{-1/2}(\partial\Omega) \hookrightarrow L^2(\partial\Omega)$ for
smooth $\partial\Omega$. $\square$
\end{proof}

\textbf{RSVP correspondence.} $\Pi_{\mathrm{IR}}$ constrains boundary values and outward
flux of $\Phi$: energy density at the surface and rate of energy exchange. It does not
directly constrain $\mathbf{v}$ or $S$ in the interior.

\subsection{LiDAR Projection}

\begin{definition}[LiDAR Projection]
Fix a threshold $\phi_0 > 0$. The LiDAR projection
$\Pi_{\mathrm{LiDAR}} : \mathcal{X} \to \mathcal{Y}_{\mathrm{LiDAR}}$ is defined by
\[
\Pi_{\mathrm{LiDAR}}(X) = \mathbf{1}_{\{\Phi \ge \phi_0\}}
\quad \text{together with} \quad
\Pi_{\mathrm{LiDAR}}^{\perp}(X) =
\frac{\nabla\Phi}{|\nabla\Phi|}\bigg|_{\{\Phi = \phi_0\}}.
\]
The observation space is
$\mathcal{Y}_{\mathrm{LiDAR}} = L^2(\Omega;\{0,1\}) \times L^2(\{\Phi=\phi_0\};S^2)$.
\end{definition}

\begin{lemma}
$\Pi_{\mathrm{LiDAR}}$ is bounded for $\phi_0$ a regular value of $\Phi$.
\end{lemma}
\begin{proof}
For $\phi_0$ a regular value, $\{\Phi = \phi_0\}$ is smooth by the implicit function
theorem in $H^1$. The indicator is bounded in $L^2(\Omega)$ by domain volume. The normal
field is bounded by the Sobolev embedding $H^1 \hookrightarrow C^0$ in three dimensions.
$\square$
\end{proof}

\textbf{RSVP correspondence.} $\Pi_{\mathrm{LiDAR}}$ constrains the geometry of level
sets of $\Phi$ and normal directions to those surfaces, which simultaneously constrain
boundary conditions on $\mathbf{v}$: flows cannot cross material boundaries.

\subsection{RGB/CCD Projection}

\begin{definition}[RGB Projection]
Let $\mathbf{c} \in S^2$ be the camera viewing direction and $f : \mathbb{R} \to \mathbb{R}^3$
a material reflectance function. The RGB projection
$\Pi_{\mathrm{RGB}} : \mathcal{X} \to \mathcal{Y}_{\mathrm{RGB}}$ is defined by
\[
\Pi_{\mathrm{RGB}}(X) = \int_0^\infty f(\Phi(\mathbf{x}(t)))\,
e^{-\int_0^t \Phi(\mathbf{x}(s)) ds}\, dt
\]
where $\mathbf{x}(t) = \mathbf{x}_0 + t\mathbf{c}$ and
$\mathcal{Y}_{\mathrm{RGB}} = L^2(\mathcal{V};\mathbb{R}^3)$ over camera pixels $\mathcal{V}$.
\end{definition}

\begin{lemma}
$\Pi_{\mathrm{RGB}}$ is bounded when $f$ is Lipschitz and $\Phi \ge 0$.
\end{lemma}
\begin{proof}
The exponential attenuation is bounded by 1. The integrand is bounded by $\|f\|_\infty$,
which is finite by Lipschitz continuity and $H^1 \hookrightarrow L^\infty$ in three
dimensions. $\square$
\end{proof}

\textbf{RSVP correspondence.} $\Pi_{\mathrm{RGB}}$ constrains $\Phi$ along rays---a
projective complement to $\Pi_{\mathrm{LiDAR}}$'s level-set constraint. Their kernels
overlap partially but not completely.

\subsection{Audio Projection}

\begin{definition}[Audio Projection]
Let $\mathbf{x}_m \in \Omega$ for $m = 1,\ldots,M$ be microphone positions. The audio
projection $\Pi_{\mathrm{audio}} : \mathcal{X} \times [0,T] \to \mathcal{Y}_{\mathrm{audio}}$
is
\[
\Pi_{\mathrm{audio}}(X)(t) =
\left( \partial_t S(\mathbf{x}_m, t) + \nabla \cdot \mathbf{v}(\mathbf{x}_m, t)
\right)_{m=1}^M
\]
with $\mathcal{Y}_{\mathrm{audio}} = L^2([0,T];\mathbb{R}^M)$.
\end{definition}

\textbf{RSVP correspondence.} $\Pi_{\mathrm{audio}}$ is the only modality that directly
constrains $\partial_t S$---the rate of entropy production---and $\nabla \cdot \mathbf{v}$,
encoding sources and sinks of flow. Hidden processes invisible to all spatial sensors are
audible as entropy production events.

\subsection{Photonic Strip Projection}

\begin{definition}[Photonic Strip Projection]
Let $\gamma : [0,L] \to \Omega$ parametrize the strip by arc length. The photonic strip
projection $\Pi_{\mathrm{strip}} : \mathcal{X} \times [0,T] \to \mathcal{Y}_{\mathrm{strip}}$
is
\[
\Pi_{\mathrm{strip}}(X)(s,t) =
\left( \Phi(\gamma(s),t),\;
\mathbf{v}(\gamma(s),t) \cdot \gamma'(s),\;
\partial_t \Phi(\gamma(s),t) \right)
\]
with $\mathcal{Y}_{\mathrm{strip}} = L^2([0,L] \times [0,T];\mathbb{R}^3)$.
\end{definition}

\textbf{RSVP correspondence.} The photonic strip provides continuous temporal records of
$\Phi$, tangential $\mathbf{v}$, and $\partial_t \Phi$ along a fixed transect: the only
modality directly constraining temporal evolution of $\Phi$ at interior points.

\subsection{The Completeness Proposition}

\begin{definition}[Spanning Condition]
A projection family $\{\Pi_i\}_{i \in \mathcal{I}}$ is \emph{field-spanning} if
\[
\ker \Ftil = \{ [X] \in \Xmod \mid \Pi_i(X) = 0 \;\forall i \in \mathcal{I} \} = \{[0]\}.
\]
\end{definition}

\begin{proposition}[Completeness of the Five-Modality Family]
\label{prop:completeness}
The projection family $\{\Pi_{\mathrm{IR}}, \Pi_{\mathrm{LiDAR}}, \Pi_{\mathrm{RGB}},
\Pi_{\mathrm{audio}}, \Pi_{\mathrm{strip}}\}$ is field-spanning on $\Xmod$.
\end{proposition}

\begin{proof}
Suppose $[X] \in \ker \Ftil$. $\Pi_{\mathrm{IR}}(X) = 0$ gives
$\Phi|_{\partial\Omega} = 0$. $\Pi_{\mathrm{LiDAR}}(X) = 0$ gives $\Phi < \phi_0$
a.e. $\Pi_{\mathrm{strip}}(X) = 0$ along strips dense in $\Omega$ gives $\Phi = 0$
on a dense subset; by $H^1$ continuity, $\Phi = 0$ on $\Omega$. With $\Phi = 0$,
strip vanishing also gives $\mathbf{v} \cdot \gamma' = 0$ along all strips; with three
non-parallel strip directions, $\mathbf{v} = 0$. Then $\Pi_{\mathrm{audio}}(X) = 0$
with $\mathbf{v} = 0$ gives $\partial_t S = 0$ at dense microphone positions, so $S$
is time-constant; RSVP entropy constraints force $S = 0$. Hence $[X] = [0]$. $\square$
\end{proof}

\subsection{Degrees of Freedom Table}

\begin{center}
\begin{tabular}{llll}
\toprule
\textbf{Modality} & \textbf{Field constrained} & \textbf{Where} &
\textbf{New DOF closed} \\
\midrule
Infrared & $\Phi$, $\nabla\Phi \cdot \hat{n}$ & $\partial\Omega$ &
Boundary energy, flux \\
LiDAR & Level sets of $\Phi$, normals & Interior surfaces &
Geometry; $\mathbf{v}$ BCs \\
RGB & $\Phi$ along rays & Lines of sight &
Interior density \\
Audio & $\partial_t S$, $\nabla \cdot \mathbf{v}$ & Point measurements &
Temporal entropy, flow sources \\
Photonic strip & $\Phi$, $\mathbf{v} \cdot \gamma'$, $\partial_t \Phi$ &
Transects over time & Temporal $\Phi$; tangential flow \\
\bottomrule
\end{tabular}
\end{center}

\subsection{Each Modality as a Chart}

Under the Sheaf-Variational Equivalence, each $\Pi_i$ defines a chart
$U_i = \{[X] \in \mathcal{A} \mid \|\Pi_i(X) - y_i\| \le \epsilon_i\}$ on $\Xmod$.
The five modalities define five charts. Completeness ensures the cover is fine enough
to support a non-degenerate sheaf. Sensor design is sheaf design: choosing modalities
is choosing the cover of $\Xmod$.

%% ============================================================
\section{The Feasible Set and Its Geometry}
\label{sec:feasible}
%% ============================================================

This section identifies the feasible set $\mathcal{F}$ as the primary geometric and
cohomological object of the Yarncrawler framework. Given the projection family
$\{\Pi_i\}_{i \in \mathcal{I}}$ and observations $\{y_i\}$, define
observation-consistent sets
\[
U_i = \left\{ [X] \in \mathcal{A} \;\middle|\;
\|\Pi_i(X) - y_i\|_{\mathcal{Y}_i} \le \epsilon_i \right\},
\qquad
\Omega_{\mathrm{obs}} = \bigcap_{i \in \mathcal{I}} U_i.
\]
The dynamically admissible set $\Gamma$, whose full construction appears in
Section~\ref{sec:disambiguation}, consists of configurations lying on admissible
trajectories of the world dynamics. The feasible set is
$\mathcal{F} = \Omega_{\mathrm{obs}} \cap \Gamma$.
Under the Sheaf-Variational Equivalence (Section~\ref{sec:equivalence}), each $U_i$
is the support of a local section of $\mathcal{S}_\Pi$, and $\mathcal{F}$ is the space
of admissible global sections. The identifiability and disambiguation theorems are
statements about when this fiber collapses to a singleton.

\subsection{Quotient Structure and Identifiability Domain}

The relevant ambient space is the quotient $\mathcal{A} = \mathcal{X}/{\sim}$, where
$X \sim X'$ if $\Pi_i(X) = \Pi_i(X')$ for all $i$. Multiple representatives in
$\mathcal{X}$ agreeing on all projections correspond to the same equivalence class in
$\mathcal{A}$ and hence to the same global section of $\mathcal{S}_\Pi$. All
identifiability statements in Section~\ref{sec:identifiability} concern the restriction
of the induced map $\Ftil : \mathcal{A} \to \prod_i \mathcal{Y}_i$ to $\mathcal{F}$.

\subsection{Geometry of the Observation-Consistent Set}

Each $U_i$ is closed and convex under Assumption~\ref{ass:regularity}. However, their
intersection $\Omega_{\mathrm{obs}}$ is generically anisotropic and nontrivial. Three
structural features determine its geometry. \emph{Anisotropy}: each projection
constrains different directions in $\mathcal{A}$. \emph{Non-aligned nullspaces}: the
kernels of $\{\Pi_i\}$ need not coincide, leaving residual directions that survive all
individual constraints. \emph{Curvature under pullback}: the metric induced by the
projection family is not Euclidean in $\mathcal{X}$, even when each $\Pi_i$ is linear.
For these reasons $\Omega_{\mathrm{obs}}$ is best understood as a stratified admissible
region whose local structure reflects the overlap pattern of the sensor family; smooth
manifold structure holds on regular strata under generic regularity assumptions.

\subsection{Dimensionality, Degeneracy, and Cohomology}

Define the effective dimension $\dim_{\mathrm{eff}}(\mathcal{F}) = \dim(\Omega_{\mathrm{obs}}
\cap \Gamma)$. When $\dim_{\mathrm{eff}}(\mathcal{F}) > 0$, multiple admissible
configurations remain; under the Sheaf-Variational Equivalence this corresponds to
multiple global sections, equivalently to non-trivial $H^1(\mathcal{U}, \mathcal{S}_\Pi)$.
Closure is the condition $\dim_{\mathrm{eff}}(\mathcal{F}) = 0$, which under the
equivalence is exactly
\[
\dim H^0(\mathcal{U}, \mathcal{S}_\Pi) = 1, \qquad H^1(\mathcal{U}, \mathcal{S}_\Pi) = 0.
\]
Ambiguity, degeneracy, and non-trivial cohomology are the same phenomenon expressed in
geometric, variational, and categorical language.

\begin{proposition}[Feasible Set Contraction]
\label{prop:contraction}
Under monotone addition of sensors with independent noise models, $\dim_{\mathrm{eff}}(\mathcal{F})$
is non-increasing. Strict decrease holds when the new sensor $\Pi_j$ satisfies
$\ker \Pi_j \not\supseteq \bigcap_{i < j} \ker \Pi_i$.
\end{proposition}
\begin{proof}
Adding $U_j$ to the intersection defining $\Omega_{\mathrm{obs}}$ cannot increase
dimension. Strict decrease holds iff $\Pi_j$ constrains a direction not previously closed.
$\square$
\end{proof}

\subsection{Information Geometry of the Feasible Set}

The projection family and noise models induce a Fisher information metric on $\mathcal{A}$:
\[
g_F(X) = \sum_{i \in \mathcal{I}} D\Pi_i(X)^* \Sigma_i^{-1} D\Pi_i(X).
\]
Directions in $\ker g_F(X)$ are observationally invisible. Two structural properties
follow. \emph{Metric monotonicity}: adding sensor $j$ gives
$g_F^{\mathrm{new}} \succeq g_F$. \emph{Dynamical lifting}: a null direction $v$ at
$t = 0$ may satisfy $g_F^{(\tau)}(X)(v,v) > 0$ under the time-extended metric, because
the flow $\Phi_F^\tau$ transports null directions into visible ones. The formal
development of both properties is in Appendix~\ref{app:fisher}.

\subsection{Information Volume and Closure}

The Fisher metric induces a volume form on the regular strata of $\mathcal{F}$. Define
the information volume $\mathrm{Vol}_F(\mathcal{F}) = \int_{\mathcal{F}_{\mathrm{reg}}}
d\mathrm{Vol}_{g_F}$. At closure the following collapse simultaneously:
\[
\mathrm{Vol}_F(\mathcal{F}) \to 0,\quad
\dim H^0 \to 1,\quad
H^1 \to 0,\quad
S_{\mathrm{RSVP}}(X^*) \text{ minimized.}
\]
These are not independent criteria but one condition in four languages.
Full equivalence is established in Appendix~\ref{app:fisher}.

\subsection{Dynamics and Non-Convex Structure}

While $\Omega_{\mathrm{obs}}$ is often convex, $\Gamma$ is generically non-convex.
Their intersection $\mathcal{F}$ may exhibit disconnected components, curved boundaries,
and trajectory-induced foliations, corresponding under the equivalence to distinct
trajectory classes and non-trivial $H^1$ classes. Section~\ref{sec:disambiguation}
shows that trajectory constraints eliminate these components by enforcing temporal
separation.

\subsection{Interpretation}

The feasible set is the space of all worlds compatible with the evidence and the laws of
motion. The task of Yarncrawler is not statistical prediction but geometric contraction:
reduce $\dim_{\mathrm{eff}}(\mathcal{F})$ until $|\mathcal{F}| = 1$. Under the
Sheaf-Variational Equivalence, this is exactly the process of repairing $\mathcal{S}_\Pi$
until a unique global section exists.

\subsection{Optimal Sensor Selection as Cohomological Descent}

The preceding subsections establish that residual ambiguity in $\mathcal{F}$ corresponds,
under the Sheaf-Variational Equivalence, to nontrivial cohomology in
$H^1(\mathcal{U}, \mathcal{S}_\Pi)$. Each unresolved degree of freedom in $\mathcal{F}$
is a generator of this cohomology group. Closure is therefore equivalent to the
annihilation of all such generators.

This allows the sensing problem to be reformulated as a descent problem on cohomological
rank. Let $r(\mathcal{U}) := \dim H^1(\mathcal{U}, \mathcal{S}_\Pi)$ and let $\phi_j$
denote the map on cohomology induced by adding sensor $j$:
$\phi_j : H^1(\mathcal{U}, \mathcal{S}_\Pi) \to H^1(\mathcal{U} \cup \{U_j\}, \mathcal{S}_\Pi)$.
The new sensor sends some cocycles to coboundaries, reducing the rank to
$r_{n+1} = \dim \mathrm{Im}(\phi_j)$.

For a candidate sensor $j$, define its marginal information gain along the residual
nullspace $N_X = \ker \Ftil|_{T_X\mathcal{F}}$ by
\[
\Delta_j(X) = \inf_{\substack{v \in N_X \\ \|v\|=1}}
\left\langle D\Pi_j(X)v,\; \Sigma_j^{-1} D\Pi_j(X)v \right\rangle.
\]

\begin{theorem}[Greedy Cohomology Descent]
\label{thm:greedy_descent}
Suppose that each sensor $j$ induces a linear action on cohomology generators via
$D\Pi_j$, the Fisher metric provides a nondegenerate quadratic form on the residual
nullspace, and sensor effects on generators are independent up to second order. Then
\[
j^* \in \arg\max_j \int_{\mathcal{F}_{\mathrm{reg}}} \Delta_j(X)\,
d\mathrm{Vol}_{g_F^{\mathcal{F}}}(X)
\;\Longleftrightarrow\;
j^* \in \arg\min_j \dim H^1(\mathcal{U} \cup \{U_j\}, \mathcal{S}_\Pi).
\]
\end{theorem}

\begin{proof}
Each generator of $H^1$ corresponds to a direction $v \in N_X$. Sensor $j$ acts on this
generator via $D\Pi_j(X)v$. If $D\Pi_j(X)v \ne 0$, the generator is detected and
eliminated: the cocycle becomes a coboundary under the refined cover. The Fisher
contribution $\langle D\Pi_j(X)v, \Sigma_j^{-1} D\Pi_j(X)v\rangle$ measures the energy
of elimination. Maximizing $\Delta_j$ integrated over $\mathcal{F}$ maximizes aggregate
elimination energy. Under independence, each generator contributes additively, so the
sensor maximizing this sum achieves maximal rank reduction in $H^1$. $\square$
\end{proof}

The criterion has a direct geometric interpretation: it selects the sensor whose
projections are most sensitive along directions that remain unconstrained by all existing
modalities. The objective is not to ``see more'' but to eliminate the specific degrees of
freedom that still admit alternative world-states.

\subsection{The Morse-Variational Structure of Yarncrawler Dynamics}

The refinement process admits a continuous limit in which the evolution of $[X]$ is
governed by a gradient flow on the consistency functional.

\begin{definition}[Yarncrawler Energy Functional]
\[
\mathcal{E}(X) = \sum_{i \in \mathcal{I}} \|\Pi_i(X) - y_i\|_{\mathcal{Y}_i}^2
+ \mathcal{R}(X).
\]
\end{definition}

The gradient flow under the Fisher metric $g_F$ is
\[
\frac{dX}{dt} = -\nabla_{g_F} \mathcal{E}(X),
\]
which decreases $\mathcal{E}$ monotonically. The critical structure of $\mathcal{E}$
recovers the cohomological content of the framework.

\begin{theorem}[Morse Structure of Residual Ambiguity]
\label{thm:morse}
Under generic conditions on $\{\Pi_i\}$ and $\mathcal{R}$, the functional $\mathcal{E}$
is a Morse function on each regular stratum of $\mathcal{F}$, and:
\begin{enumerate}
\item[(i)] Critical points of $\mathcal{E}$ correspond to equivalence classes locally
indistinguishable under all projections.
\item[(ii)] The index of a critical point equals $\dim N_X$, the number of independent
ambiguity directions.
\item[(iii)] Nontrivial cohomology classes in $H^1(\mathcal{U}, \mathcal{S}_\Pi)$
correspond to non-minimal critical points of $\mathcal{E}$.
\item[(iv)] The unique global minimum corresponds to the closure state $[X^*]$ with
$\mathcal{E}(X^*) = 0$.
\end{enumerate}
\end{theorem}

\begin{proof}[Proof sketch]
At a critical point $X_c$, $\nabla\mathcal{E}(X_c) = 0$, which holds when all projection
residuals vanish or the gradient is in equilibrium with $\nabla\mathcal{R}$. The Hessian
$\nabla^2\mathcal{E}(X_c) = g_F(X_c) + \nabla^2\mathcal{R}(X_c)$ has nullspace equal to
directions invisible to all projections and unpenalized by $\mathcal{R}$---exactly
the generators of $H^1$. Non-minimal critical points have positive-index nullspace
corresponding to these generators. $\square$
\end{proof}

\begin{corollary}[Closure as Unique Morse Minimum]
If the conditions of Theorems~\ref{thm:identifiability} and~\ref{thm:disambiguation}
hold, then $\mathcal{E}$ has a unique global minimum $[X^*]$ and all other critical points
are unstable. Equivalently:
\[
\text{all cohomology classes trivialized}
\;\Longleftrightarrow\;
\mathcal{E} \text{ has a unique minimum}.
\]
\end{corollary}

\subsection{Distributional Lift: Wasserstein Gradient Flow}

The pointwise framework admits a distributional generalization. Treat the state not as a
point estimate but as a probability measure $\mu \in \mathcal{P}(\mathcal{A})$ supported
on the current feasible set.

\begin{definition}[Measure-Valued Energy Functional]
\[
\mathcal{E}[\mu] = \int_{\mathcal{A}} \left(
\sum_i \|\Pi_i(X) - y_i\|^2 + \mathcal{R}(X)
\right) d\mu(X) + \tau\, \mathrm{Ent}(\mu),
\]
where $\mathrm{Ent}(\mu) = \int \mu \log \mu$ is the Boltzmann entropy and $\tau > 0$
is a temperature parameter encoding exploration-exploitation balance.
\end{definition}

Endowing $\mathcal{P}(\mathcal{A})$ with the 2-Wasserstein metric $W_2$, the gradient
flow of $\mathcal{E}[\mu]$ is the continuity equation
\[
\partial_t \mu_t = \nabla \cdot \left(
\mu_t \nabla \frac{\delta \mathcal{E}}{\delta \mu}
\right),
\qquad
\frac{\delta \mathcal{E}}{\delta \mu}(X) =
\sum_i \|\Pi_i(X) - y_i\|^2 + \mathcal{R}(X) + \tau \log \mu_t(X).
\]
Under the conditions of Theorems~\ref{thm:identifiability} and~\ref{thm:disambiguation},
as $t \to \infty$ and $\tau \to 0$:
\[
\mu_t \;\longrightarrow\; \delta_{[X^*]}, \qquad \mathrm{supp}(\mu_t) \to \{[X^*]\}.
\]
The pointwise consistency operator is recovered as the zero-temperature limit:
$\tau \to 0 \Rightarrow \mu_t \to \delta_{X(t)}$, recovering
$\frac{dX}{dt} = -\nabla\mathcal{E}(X)$.

In the distributional formulation, cohomological obstruction manifests as multi-modality
of $\mu_t$: distinct connected components of $\mathrm{supp}(\mu_t)$ correspond to
distinct global sections. The elimination of cohomology is the transition from
multi-modal to unimodal concentration.

\subsection{RSVP Realization of the Variational Flow}

The Wasserstein gradient flow of Section~\ref{sec:feasible}.8 is identical to RSVP
field evolution under a constraint-induced potential. This establishes that the Yarncrawler
refinement process is not merely analogous to RSVP dynamics; it is a specific
instantiation of RSVP dynamics under constraint-driven forcing.

Represent the measure $\mu_t$ as a density field: $\mu_t(X) \leftrightarrow \Phi(x,t)$.
The continuity equation $\partial_t \mu_t + \nabla\cdot(\mu_t u_t) = 0$ becomes the
RSVP scalar conservation law $\partial_t \Phi + \nabla\cdot(\Phi\mathbf{v}) = 0$,
with velocity field
\[
\mathbf{v} = -\nabla V_{\mathrm{proj}} - \nabla V_{\mathrm{adm}} + \tau \nabla S,
\]
where $V_{\mathrm{proj}}(X) = \sum_i \|\Pi_i(X) - y_i\|^2$ is the projection potential,
$V_{\mathrm{adm}}(X) = \mathcal{R}(X)$ is the admissibility potential, and
$S = -\log\Phi$ is the RSVP entropy field (so $-\tau\nabla\log\Phi = \tau\nabla S$).

The resulting field equation is
\[
\partial_t \Phi + \nabla\cdot\!\left[
\Phi\!\left(-\nabla V_{\mathrm{proj}} - \nabla V_{\mathrm{adm}} + \tau\nabla S\right)
\right] = 0,
\]
a forced advection-diffusion equation in which projection and admissibility terms drive
directed collapse toward constraint-satisfying configurations, while the entropy term
induces diffusive spreading preventing premature concentration.

Closure occurs when: (i) projection mismatch vanishes, $\Pi_i(X^*) = y_i$; (ii)
admissibility is satisfied, $\mathcal{R}(X^*) = \min$; and (iii) entropy collapses,
$S \to -\infty$ ($\Phi \to \delta_{[X^*]}$), giving $\partial_t\Phi \to 0$.

The full equivalence is:
\[
\text{Yarncrawler refinement}
\;\equiv\;
\text{RSVP field evolution}
\;\equiv\;
\text{Wasserstein gradient flow.}
\]
These are not three descriptions. They are one dynamical system written in categorical
form (sheaf repair), geometric form (feasible set collapse), variational form (energy
minimization), and physical form (RSVP fields). Every layer of the framework is the same
object: sheaf theory eliminates $H^1$, geometry collapses $\dim\mathcal{F}$, information
theory collapses $\mathrm{Vol}_F(\mathcal{F})$, dynamics execute the gradient flow of
$\mathcal{E}$, and optimization drives the convergence of $\mathcal{C}$. These are
different coordinate systems on one process.

%% ============================================================
\section{The Consistency Operator and the Identifiability Theorem}
\label{sec:identifiability}
%% ============================================================

\subsection{Setup and Standing Assumptions}

Let $\mathcal{X} = H^1(\Omega) \times H^1(\Omega;\mathbb{R}^3) \times L^2(\Omega)$
with the product Sobolev topology (Appendix~\ref{app:frechet}). A symmetry group $G$
acts smoothly, properly, and isometrically on $\mathcal{X}$; the quotient $\Xmod$ is
Hausdorff. We work in $\Xmod$ throughout.

\begin{assumption}[Projection Regularity]
\label{ass:regularity}
Each $\Pi_i : \mathcal{X} \to \mathcal{Y}_i$ is a bounded linear map to a separable
Hilbert space $\mathcal{Y}_i$.
\end{assumption}

\begin{assumption}[Noise Model]
\label{ass:noise}
Each observation satisfies $y_i = \Pi_i(X^*) + \eta_i$ where $\|\eta_i\|_{\mathcal{Y}_i}
\le \epsilon_i$ almost surely.
\end{assumption}

\begin{assumption}[Regularizer Structure]
\label{ass:regularizer}
The regularizer takes the form
\[
\mathcal{R}(X) = \mathcal{R}_{\mathrm{RSVP}}(X)
+ \lambda_{\mathrm{eco}}\,\mathcal{R}_{\mathrm{eco}}(X)
+ \lambda_{\mathrm{civ}}\,\mathcal{R}_{\mathrm{civ}}(X)
\]
where $\mathcal{R}_{\mathrm{RSVP}}$ measures violation of the RSVP field equations in
$L^2$, $\mathcal{R}_{\mathrm{eco}}$ is zero on ecologically admissible configurations
and strictly positive off them, and $\mathcal{R}_{\mathrm{civ}}$ is zero on
civilizationally admissible configurations and strictly positive off them.
\end{assumption}

\begin{assumption}[Admissibility Region]
\label{ass:admissibility}
The admissible set $\mathcal{A} = \{[X] \in \Xmod \mid \mathcal{R}(X) < \infty\}$
is a closed convex subset of $\Xmod$, and $\mathcal{R}$ restricted to $\mathcal{A}$
is strictly convex along geodesics.
\end{assumption}

\begin{assumption}[Banach Realization]
\label{ass:banach}
There exists a Banach space $\mathcal{B}$ and continuous embedding
$\iota : \mathcal{X} \hookrightarrow \mathcal{B}$ such that $\mathcal{A}$ is closed and
convex in $\mathcal{B}$, each $\Pi_i$ extends to a bounded linear operator on
$\mathcal{B}$, and $\mathcal{R}$ extends to a proper lower semicontinuous functional on
$\mathcal{B}$.
\end{assumption}

In practice one takes $\mathcal{B} = \mathcal{X}$ itself, which is already a Hilbert
space. The Banach realization is stated separately to make the compactness and weak
convergence arguments in Appendix~\ref{app:identifiability} self-contained.

\subsection{The Feasible Set}

\begin{definition}[Observation-Consistent Set]
\[
\Omega_{\mathrm{obs}} = \{ [X] \in \mathcal{A} \mid
\|\Pi_i(X) - y_i\|_{\mathcal{Y}_i} \le \epsilon_i \;\; \forall i \in \mathcal{I} \}.
\]
\end{definition}

\begin{lemma}[Non-emptiness]
Under Assumptions~\ref{ass:regularity} and \ref{ass:noise}, $X^* \in \Omega_{\mathrm{obs}}$
almost surely.
\end{lemma}
\begin{proof}
$\|y_i - \Pi_i(X^*)\| = \|\eta_i\| \le \epsilon_i$ almost surely by Assumption~\ref{ass:noise}.
$\square$
\end{proof}

\begin{lemma}[Closedness and Convexity]
Under Assumptions~\ref{ass:regularity} and \ref{ass:admissibility},
$\Omega_{\mathrm{obs}}$ is a closed convex subset of $\mathcal{A}$.
\end{lemma}
\begin{proof}
Each constraint set is closed (continuity of $\Pi_i$) and convex (linearity of $\Pi_i$,
convexity of norm ball). Intersection of closed convex sets is closed and convex.
$\square$
\end{proof}

The dynamically admissible set $\Gamma$ is introduced as a named object here; its full
definition and construction are given in Section~\ref{sec:disambiguation}. The
\emph{feasible set} is $\mathcal{F} = \Omega_{\mathrm{obs}} \cap \Gamma$.

\subsection{The Consistency Operator}

\begin{definition}[Consistency Operator]
\[
\mathcal{C}(X) = \arg\min_{[X'] \in \mathcal{A}}
\left[ \sum_{i \in \mathcal{I}} \|\Pi_i(X') - y_i\|_{\mathcal{Y}_i}^2
+ \mathcal{R}(X') \right].
\]
\end{definition}

$\mathcal{C}$ maps $\mathcal{A}$ to $\mathcal{A}$: the optimization is constrained to the
admissible region. A configuration minimizing projection error but violating RSVP equations
or ecological plausibility is not a candidate world-state. The regularizer defines the
domain of the operator, not a preference over an unconstrained space.

\begin{lemma}[Well-definedness]
Under Assumptions~\ref{ass:regularity}--\ref{ass:banach}, $\mathcal{C}(X)$ exists and is
unique for every $[X] \in \mathcal{A}$.
\end{lemma}
\begin{proof}
The objective $f(X') = \sum_i \|\Pi_i(X') - y_i\|^2 + \mathcal{R}(X')$ is continuous,
strictly convex (sum of convex functions plus strictly convex $\mathcal{R}$), and coercive
on $\mathcal{A}$. A continuous strictly convex coercive function on a closed convex set
attains its minimum at a unique point. Full proof in Appendix~\ref{app:identifiability}.
$\square$
\end{proof}

\subsection{The Identifiability Theorem}

The induced map of the projection family on the quotient is:
\[
\Ftil : \Xmod \longrightarrow \prod_{i \in \mathcal{I}} \mathcal{Y}_i, \qquad
[X] \mapsto (\Pi_i(X))_{i \in \mathcal{I}}.
\]
By definition of $\sim$, $\Ftil$ is well-defined.

\begin{theorem}[Yarncrawler Identifiability]
\label{thm:identifiability}
Under Assumptions~\ref{ass:regularity}--\ref{ass:banach}, the consistency operator
$\mathcal{C}$ has a unique fixed point $[X^*] \in \mathcal{F}$ if and only if $\Ftil$
is injective on $\mathcal{F}$.
\end{theorem}

\begin{proof}
$(\Rightarrow)$ Suppose $\mathcal{C}$ has a unique fixed point $[X^*] \in \mathcal{F}$.
If $\Ftil$ were not injective on $\mathcal{F}$, there exist distinct $[X],[X'] \in \mathcal{F}$
with $\Pi_i(X) = \Pi_i(X')$ for all $i$. Both achieve zero projection residual in the
objective. The objective reduces to $\mathcal{R}(X)$ vs.\ $\mathcal{R}(X')$. Either
$\mathcal{R}(X) = \mathcal{R}(X')$---violating strict convexity of $\mathcal{R}$ on
$\mathcal{A}$ (Assumption~\ref{ass:admissibility})---or one has strictly smaller
regularizer value, making the other not a minimizer. Either way contradicts uniqueness.

$(\Leftarrow)$ If $\Ftil$ is injective on $\mathcal{F}$, then $\Omega_{\mathrm{obs}} \cap
\mathcal{F}$ contains at most one point. By Lemma 5.1, $[X^*] \in \Omega_{\mathrm{obs}}$;
by definition $[X^*] \in \Gamma$; hence $[X^*] \in \mathcal{F}$. Injectivity forces
$|\mathcal{F}| = 1$. At $[X^*]$ all projection residuals are within tolerance and
$\mathcal{R}(X^*)$ is minimized over the singleton $\mathcal{F}$. Hence
$\mathcal{C}([X^*]) = [X^*]$. $\square$
\end{proof}

\subsection{The Corollary on Diagnosable Failure Modes}

\begin{corollary}
\label{cor:failures}
Identifiability fails in exactly two geometrically distinguishable ways:

\emph{(i) Projection degeneracy}: $\Ftil$ is not injective on $\mathcal{F}$. Two distinct
admissible world-states produce identical observations. This failure is intrinsic to the
sensor family and cannot be resolved by the regularizer.

\emph{(ii) Regularizer flatness}: $\mathcal{R}$ is not strictly convex on $\mathcal{F}$.
The regularizer assigns equal admissibility to configurations the sensors cannot
distinguish. This failure is intrinsic to the model class.

These failure modes are geometrically distinguishable: projection degeneracy corresponds
to a non-trivial kernel of $\Ftil|_{\mathcal{F}}$; regularizer flatness corresponds to a
flat direction of $\mathcal{R}$ within $\mathcal{F}$. They can coexist.
\end{corollary}

\subsection{Why the Quotient Structure Is Not Optional}

In a standard inverse problem one works in $\mathcal{X}$ directly. But in the Yarncrawler
setting, two configurations related by a gauge transformation or coordinate
reparametrization represent the same physical situation. Working in $\mathcal{X}$ rather
than $\Xmod$ would make uniqueness false by construction---there would always be a
continuum of configurations achieving the same objective, related by the action of $G$.

An $L^2$ uniqueness statement conflates field-theoretic gauge redundancy with measurement
noise. A measure-theoretic statement collapses the geometry of $\Xmod$ into observational
statistics and loses the regularizer's role as an admissibility condition. The
symmetry-group formulation is not a choice among equivalent formalisms; it is the only
formulation that preserves the semantic commitments of the framework.

%% ============================================================
\section{Trajectory Learning and the Dynamical Admissibility Set}
\label{sec:disambiguation}
%% ============================================================

\subsection{The Central Obstruction}

Section~\ref{sec:identifiability} established that $\mathcal{C}$ has a unique fixed point
iff $\Ftil$ is injective on $\mathcal{F}$. This condition is clean but generically unmet.
Under realistic sensing, $\Omega_{\mathrm{obs}}$ is high-dimensional. Projection degeneracy
is the default condition, not a pathological case.

The resolution cannot come from within the static framework. What is required is a
constraint of a different kind---one operating not on the configuration $X$ at a single
time but on the trajectory through which $X$ was reached. The dynamical admissibility set
$\Gamma$ is not an additional observation; it is a restriction on which configurations are
reachable under the dynamics of the world. Two configurations in $\Omega_{\mathrm{obs}}$
that are observationally indistinguishable at one moment may evolve under those dynamics
into configurations that are distinguishable at a later moment.

\subsection{The Path Space and the Stochastic Kernel}

\begin{definition}[Discrete Path Space]
$\mathcal{P}_\Delta = \{ \mathbf{X} = ([X_0],[X_1],\ldots,[X_N]) \mid [X_k] \in \mathcal{A},\;
N = \lfloor T/\Delta \rfloor \}$.
\end{definition}

\begin{definition}[Stochastic Transition Kernel]
$T_\Delta : \mathcal{A} \to \mathcal{P}(\mathcal{A})$, where for each $[X] \in \mathcal{A}$,
$T_\Delta([X],\cdot)$ specifies the distribution over one-step successors under the world
dynamics at resolution $\Delta$.
\end{definition}

\begin{definition}[Dynamically Admissible Set]
\[
\Gamma = \left\{ [X] \in \mathcal{A} \;\middle|\;
\exists\, [X_0] \in \mathcal{A},\; \exists\, k \ge 0 \text{ s.t. }
\mu_{X_0}^{T_\Delta}([X_k] = [X]) > 0 \right\}.
\]
$\Gamma$ is the support of the union of all admissible path measures. It is a geometric
object---a subset of $\mathcal{A}$---not a probabilistic one.
\end{definition}

\begin{lemma}[Invariance of $\Gamma$]
If $T_\Delta$ is time-homogeneous and $\mathcal{A}$ is forward-invariant under $T_\Delta$,
then $\Gamma \subseteq \mathcal{A}$ and $\Gamma$ is forward-invariant.
\end{lemma}

\subsection{The Deterministic Limit and Tangent Bundle Constraint}

\begin{assumption}[Generator Regularity]
\label{ass:generator}
There exists a Lipschitz vector field $F(\cdot\,;\theta) : \mathcal{A} \to T\mathcal{A}$
such that $T_\Delta([X],\cdot) \Rightarrow \delta_{[X]+\Delta F([X];\theta)}$ weakly as
$\Delta \to 0$.
\end{assumption}

In the zero-noise limit, the dynamically admissible set reduces to:
\[
\Gamma^0 = \left\{ [X] \in \mathcal{A} \;\middle|\;
\exists\, [X_0] \in \mathcal{A},\; \exists\, t \in [0,T] \text{ s.t. }
\Phi^t_F([X_0]) = [X] \right\}
\]
where $\Phi^t_F$ is the flow of $F$ at time $t$.

\begin{definition}[Admissible Tangent Cone]
For $[X] \in \Gamma^0$:
$\mathcal{T}_{[X]}\Gamma^0 = \{v \in T_{[X]}\mathcal{A} \mid v = F([X];\theta)
\text{ for some } \theta \in \Theta\}$.
\end{definition}

\begin{proposition}[Tangent Dimension Reduction]
If $\Theta \subset \mathbb{R}^p$, then $\dim \mathcal{T}_{[X]}\Gamma^0 \le p$.
In particular, if $p < \dim T_{[X]}\mathcal{A}$, the dynamical constraint reduces the
effective dimensionality of admissible directions from $\dim T_{[X]}\mathcal{A}$ to
at most $p$.
\end{proposition}
\begin{proof}
The map $\theta \mapsto F([X];\theta)$ is smooth from $\Theta \subset \mathbb{R}^p$ to
$T_{[X]}\mathcal{A}$; its image has dimension at most $p$ by the rank theorem. $\square$
\end{proof}

\subsection{Trajectory Learning from Before/After Pairs}

\begin{definition}[Fragment-Consistent Parameter Set]
Given fragment $(y_{\mathrm{before}}, y_{\mathrm{after}}, \tau)$:
\[
\Theta_{(y_b,y_a,\tau)} = \left\{ \theta \in \Theta \;\middle|\;
\exists\, [X_b] \in \Omega_{\mathrm{obs}}(y_b),\;
\Pi_i(\Phi^\tau_F([X_b])) \in B_{\epsilon_i}(y_{a,i})\; \forall i \right\}.
\]
Given training set $\mathcal{D}$, the learned parameter set is
$\hat{\Theta} = \bigcap_{n} \Theta_{(y_b^{(n)}, y_a^{(n)}, \tau^{(n)})}$.
\end{definition}

\textbf{The operator class library.} In practice $F(\cdot\,;\theta)$ is decomposed as:
\[
F([X];\theta) = \sum_{k=1}^K \alpha_k([X]) \cdot F_k([X];\theta_k)
\]
where each $F_k$ corresponds to a named ecological or civilizational process---forest
succession, urban expansion, agricultural cycling, infrastructure decay, hydrological
change, demographic transition. Each $F_k$ is learned from before/after pairs in which
that process is dominant. The activation weights $\alpha_k([X]) \ge 0$ satisfy
$\sum_k \alpha_k = 1$.

\subsection{The Dynamical Disambiguation Theorem}

\begin{definition}[Trajectory Separation]
Two configurations $[X],[X'] \in \Omega_{\mathrm{obs}}$ are \emph{trajectory-separated} by
$F(\cdot\,;\hat{\theta})$ at time $\tau > 0$ if
\[
\left\| \Ftil(\Phi^\tau_F([X];\hat{\theta})) - \Ftil(\Phi^\tau_F([X'];\hat{\theta}))
\right\|_{\prod_i \mathcal{Y}_i} > \epsilon
\]
for some $\epsilon > 0$ depending only on the noise levels $\{\epsilon_i\}$.
\end{definition}

\begin{theorem}[Dynamical Disambiguation]
\label{thm:disambiguation}
Under Assumptions~\ref{ass:regularity}--\ref{ass:generator}, if $|\Omega_{\mathrm{obs}}| > 1$
and every pair $[X],[X'] \in \Omega_{\mathrm{obs}}$ with $[X] \ne [X']$ is
trajectory-separated by $F(\cdot\,;\hat{\theta})$ at some $\tau \in (0,T]$, then
$|\mathcal{F}| = |\Omega_{\mathrm{obs}} \cap \Gamma^0| = 1$.
\end{theorem}

\begin{proof}
Let $[X],[X'] \in \Omega_{\mathrm{obs}} \cap \Gamma^0$ be distinct. Since both lie in
$\Gamma^0$, both lie on integral curves of $F(\cdot\,;\hat{\theta})$. By trajectory
separation, there exists $\tau \in (0,T]$ such that the forward-evolved projections differ
by more than $\epsilon$. At time $\tau$, observed data $y(\tau)$ must lie within $\epsilon_i$
of the true trajectory; therefore at most one of $[X(\tau)],[X'(\tau)]$ lies in
$\Omega_{\mathrm{obs}}(y(\tau))$. The true state $[X^*(\tau)] \in \Omega_{\mathrm{obs}}(y(\tau))$
by Lemma 5.1 at time $\tau$. Applying to all pairs eliminates all but $[X^*]$.
Hence $|\mathcal{F}| = 1$. $\square$
\end{proof}

\subsection{Joint Identifiability}

\begin{corollary}[Joint Identifiability]
The consistency operator $\mathcal{C}$ restricted to $\mathcal{F}$ has a unique fixed point
if either: (i) $\Ftil$ is injective on $\mathcal{F}$; (ii) every pair in
$\Omega_{\mathrm{obs}} \cap \Gamma^0$ is trajectory-separated; or (iii) both conditions
hold partially. The sufficient condition for (iii) is that $\dim(\Omega_{\mathrm{obs}}/\ker\Ftil)
+ \dim(\Gamma^0/\Gamma^0_{\mathrm{indist}}) \ge \dim\Omega_{\mathrm{obs}}$.
\end{corollary}

Case (iii) is the generic realistic scenario: sensors eliminate most ambiguity, trajectory
learning eliminates the residual. The two mechanisms are additive in their dimensional
reduction, not competing.

%% ============================================================
\section{The Inversion Layer}
\label{sec:inversion}
%% ============================================================

\subsection{Role Specification}

The inversion layer is not the model. It is a proposal mechanism: its role is to place
the iterative optimization in the basin of the correct fixed point. The consistency
operator $\mathcal{C}$ is the model; the encoder is a high-entropy initializer; the
physics-constrained forward model is the contraction operator. Conflating encoder with
solution is the failure mode of standard encoder-decoder stacks: doing so makes the two
failure modes of Corollary~\ref{cor:failures} undiagnosable. Separating proposal from
decision preserves the falsifiability structure.

\subsection{The Encoder as Approximate Posterior}

The encoder implements an approximate posterior:
\[
q_\phi(X \mid \{y_i\}) \approx p(X \mid \{y_i\}) \propto
\exp\!\left(-\sum_i d_i(\Pi_i(X), y_i) - \mathcal{R}(X)\right).
\]
The encoder approximates this distribution; it does not solve it. The forward-consistency
loop is what contracts that distribution. Because $\mathcal{R}$ encodes physical
admissibility, the approximate posterior is already biased toward physically plausible
configurations: the encoder does not need to learn the RSVP dynamics from scratch, only
to propose configurations that are approximately consistent with them. The encoder
minimizes KL divergence from the true posterior, subject to its architectural constraints,
treating $\mathcal{R}$ as an implicit prior.

\subsection{The Refinement Loop}

The five-step loop is:
\begin{enumerate}
\item Encoder samples $X^{(0)} \sim q_\phi(\cdot \mid \{y_i\})$
\item Forward model computes $\{\hat{y}_i\} = \{\Pi_i(X^{(0)})\}$ using the operators
  of Section~\ref{sec:projections}
\item Residuals $r_i = d_i(\hat{y}_i, y_i)$ drive a projected gradient update
\item $X^{(k+1)} = \mathcal{C}(X^{(k)})$, one step of projected gradient descent on
  the consistency functional
\item Iterate until $\mathrm{Vol}_F(\mathcal{F}) < \varepsilon$ (closure criterion of
  Section~\ref{sec:closure})
\end{enumerate}
Theorem~\ref{thm:identifiability} guarantees a unique fixed point when the identifiability
conditions hold. The refinement loop approximates the path to that fixed point. Convergence
rate is governed by the condition number of $\nabla^2\mathcal{E}$ at $[X^*]$, equal to
$g_F(X^*) + \nabla^2\mathcal{R}(X^*)$, as established in Appendix~\ref{app:fisher}.

\subsection{Categorical Structure of Multi-Modal Fusion}

The encoder decomposes into modality-specific proposal maps
$q_\phi^{(i)} : \mathcal{Y}_i \to \mathcal{P}(\mathcal{A})$ composed via a lax monoidal
structure over $\prod_i \mathcal{Y}_i$. Each modality-specific encoder produces a marginal
proposal; the fusion map combines them before the consistency operator acts. This describes
how multi-modal proposals are combined, but the explanatory weight rests on the posterior
formulation, not on the categorical structure. The categorical framing is reserved for
structural analysis of how proposals compose across deployment scales
(Section~\ref{sec:deployment}).

\subsection{The Forward Model}

Given a proposed $[X]$, the forward model produces $\{\Pi_i(X)\}$ using the operators of
Section~\ref{sec:projections}. For infrared and LiDAR, this is direct evaluation of the
trace and level-set operators. For RGB, it is a volume rendering pass using the NeRF
integral. For audio, it requires integrating the RSVP field equations forward in time from
$[X]$ to produce $\partial_t S$ and $\nabla\cdot\mathbf{v}$ at microphone positions. For
photonic strips, it evaluates the transect trace operator over $[0,T]$. The computational
cost is highest for audio and strip projections; learned surrogate models replace the full
forward pass in those cases during iterative refinement, with the exact operators used only
for final validation.

\subsection{Why Not End-to-End}

An end-to-end system collapses proposal and decision into a single learned map
$\{y_i\} \mapsto [X]$. This replaces $\mathcal{C}$ with the encoder and discards the
forward model. The result cannot diagnose its own failure modes (Corollary~\ref{cor:failures}
becomes inaccessible), cannot incorporate new sensor modalities without retraining, and
cannot be falsified by physical constraints. The separation of encoder from consistency
operator is not an engineering preference. It is required by the identifiability
structure: semantics are not learned directly; they are selected via constraint closure.

%% ============================================================
\section{Closure: Definition, Criteria, and Detection}
\label{sec:closure}
%% ============================================================

\subsection{The Unified Definition}

A state $[X^*] \in \Xmod$ is \emph{closed} if it is the unique fixed point of
$\mathcal{C}$ and lies on a trajectory of $T_\Delta$ consistent with all observed
fragments. This definition decomposes into three operationally distinct conditions that
together constitute closure in a deployed system.

\subsection{The Three Conditions}

\emph{Projection parity}: $\|\Pi_i(X^*) - y_i\|_{\mathcal{Y}_i} \le \epsilon_i$ for
all $i \in \mathcal{I}$. The fixed point is consistent with all current observations.
This is verifiable in real time by evaluating the forward model.

\emph{Dynamical stability}: $[X^*] \in \Gamma^0$, meaning $[X^*]$ lies on an integral
curve of $F(\cdot\,;\hat{\theta})$. The fixed point is reachable under the learned
dynamics, not merely consistent with instantaneous observations. This rules out
configurations that fit current data but could not have been produced by any admissible
trajectory.

\emph{Counterfactual consistency}: if $[X^*]$ is simulated forward under $F$ for time
$\tau$ and reprojected, the result must lie within $\epsilon_i$ of actual observations
at time $\tau$. This is the strongest condition: it requires the fixed point's forward
trajectory to be consistent with future observations, not just its current projection.

\subsection{Equivalence of the Three Conditions}

\begin{theorem}[Equivalence of Closure Criteria]
\label{thm:closure_equiv}
Under the assumptions of Theorems~\ref{thm:identifiability} and~\ref{thm:disambiguation},
the three conditions of Section~8.2 are equivalent. Each implies
$\dim_{\mathrm{eff}}(\mathcal{F}) = 0$, and $\dim_{\mathrm{eff}}(\mathcal{F}) = 0$
implies all three.
\end{theorem}

\begin{proof}
Projection parity tightens $\Omega_{\mathrm{obs}}$ until $[X^*]$ is its only element.
Dynamical stability ensures $[X^*] \in \Gamma^0$, so $\mathcal{F} = \{[X^*]\}$. Under
the Sheaf-Variational Equivalence, $\dim_{\mathrm{eff}}(\mathcal{F}) = 0$ is equivalent
to $H^0 = 1$ and $H^1 = 0$. Counterfactual consistency is the temporal extension of
projection parity: the sheaf condition holds over the spatio-temporal cover, not just the
spatial cover. Each condition therefore implies $|\mathcal{F}| = 1$, and $|\mathcal{F}|
= 1$ implies all three. $\square$
\end{proof}

The three conditions are not independent criteria to be checked separately. They are
three faces of the single condition $|\mathcal{F}| = 1$.

\subsection{Operational Detection}

Closure is operationally detectable as stability of $[X^*]$ under perturbation of
$\epsilon_i$: if tightening any constraint by $\delta\epsilon_i$ does not move the fixed
point, closure has been achieved. This provides the stopping criterion for the refinement
loop of Section~\ref{sec:inversion}. The rate of convergence to closure, as a function of
sensor density and trajectory fragment coverage, is governed by the eigenspectrum of
$g_F(X^*) + \nabla^2\mathcal{R}(X^*)$ established in Appendix~\ref{app:fisher}.

\subsection{Entropy Collapse at Closure}

At closure, the following hold simultaneously:
$\mathrm{Vol}_F(\mathcal{F}) = 0$, $H^1(\mathcal{U}, \mathcal{S}_\Pi) = 0$,
$\dim H^0 = 1$, and $S_{\mathrm{RSVP}}(X^*)$ is minimized subject to admissibility.
By Theorem~\ref{thm:closure_volume} and Corollary~\ref{cor:closure_criterion}, these are
equivalent under nondegeneracy of $g_F^{\mathcal{F}}$. The RSVP entropy field at closure
is the local density of eliminated cohomological obstruction. The closed state is not a
best estimate; it is the unique trajectory consistent with all projections and admissible
dynamics.

\subsection{Failure Modes Revisited}

Corollary~\ref{cor:failures} identified two failure modes. Projection degeneracy
($H^1 \ne 0$) means closure cannot be achieved by additional observation at the current
moment: it requires new sensor modalities (refining the cover) or temporal extension
(Section~\ref{sec:disambiguation}'s disambiguation mechanism). Regularizer flatness means
closure cannot be achieved by any sensor: it requires revision of the model class
$\mathcal{A}$ itself. These are the only two ways the refinement loop can fail to
converge, and both are diagnosable by the criteria of this section.

%% ============================================================
\section{Deployment Scales as Feasible Set Geometry}
\label{sec:deployment}
%% ============================================================

\subsection{The Unified Framework}

All three deployment scales are instances of the same consistency operator $\mathcal{C}$
with different $(\Pi_i, \mathcal{N}_i, \epsilon_i)$ profiles. The scale is not a
separate architectural choice; it is a parameter of the noise model and the available
projection family. The identifiability and disambiguation theorems apply uniformly across
scales.

\subsection{Planetary Scale}

At planetary scale, sensors are satellite multispectral imagery (RGB and infrared) with
occasional LiDAR passes; audio and photonic strips are absent. Temporal resolution is
weeks to months; spatial resolution 10--100 meters per pixel; noise levels $\epsilon_i$
are large. The feasible set $\Omega_{\mathrm{obs}}$ is high-dimensional and weakly
constrained. The system operates in the high-entropy proposal regime, relying heavily on
$\mathcal{R}$ to restrict $\mathcal{F}$. Trajectory learning carries the disambiguation
burden: seasonal cycles, annual land cover change, and multi-year succession patterns
provide the operator class library entries that contract $\Gamma$. The projection family
is incomplete (Proposition~\ref{prop:completeness} fails due to missing audio and strip
modalities), but Corollary~\ref{cor:joint_identifiability} applies: trajectory learning
compensates for the missing static constraints.

\subsection{Regional Scale}

At regional scale, sensors include aerial multispectral and LiDAR surveys, distributed
acoustic sensors, and periodic photonic strip transects. Temporal resolution is days to
weeks; spatial resolution 1--10 meters. $\Omega_{\mathrm{obs}}$ contracts along spatial
dimensions but remains temporally undersampled. Both static projection richness and
trajectory learning are required. The transition from planetary to regional scale is a
cover refinement: the planetary-scale sheaf $\mathcal{S}_\Pi$ restricts to the
regional-scale sheaf $\mathcal{S}_{\Pi'}$ on the finer cover, with $H^1$ classes killed
by the additional modalities.

\subsection{Site Scale}

At site scale, fixed arrays include all five modalities at high spatial and temporal
density: seconds to minutes temporal resolution, centimeters to meters spatial resolution,
small $\epsilon_i$. $\Omega_{\mathrm{obs}}$ collapses rapidly toward a singleton. The
system achieves genuine closure. Re-projection residuals (forward model outputs against
actual sensor readings) provide a continuous closure detection signal implementing the
operational criterion of Section~\ref{sec:closure}.

\subsection{Multi-Scale Fusion and Scale Consistency}

\begin{theorem}[Scale Consistency]
\label{thm:scale}
If the noise models $\{\mathcal{N}_i\}$ are correctly specified at each scale, the fixed
point $[X^*]$ is invariant to which scale drives the final refinement step.
\end{theorem}
\begin{proof}
The functor $\Psi$ of Theorem~\ref{thm:equivalence} maps cover refinements to sheaf
restrictions. The restriction map from the finer (site-scale) sheaf to the coarser
(planetary-scale) sheaf preserves the global section, since the global section at the
finer scale restricts to a global section at the coarser scale. Uniqueness at each scale
(under the identifiability conditions) forces these to agree. $\square$
\end{proof}

In practice: planetary constraints provide the prior on $\mathcal{R}$, regional
constraints initialize $q_\phi$, and site constraints drive the refinement loop to
closure. The three scales are nested constraint systems, not separate deployments.

\subsection{Sensor Selection Across Scales}

The optimal sensor selection criterion of Theorem~\ref{thm:greedy_descent} applies
uniformly across scales. At planetary scale, the most valuable next sensor is the one
that closes the largest-dimensional direction in $\ker \Ftil|_{\mathcal{F}}$: typically
audio (constrains $\partial_t S$ and $\nabla\cdot\mathbf{v}$, both unconstrained by
satellite imagery) or high-frequency temporal sampling (constrains $\partial_t\Phi$).
At site scale, where most static degrees of freedom are already closed, the criterion
favors temporal extension of the observation window over additional simultaneous modalities.

%% ============================================================
\section{Relation to Prior Work}
\label{sec:related}
%% ============================================================

\subsection{Data Assimilation and Ensemble Methods}

Data assimilation methods, including ensemble Kalman filtering (EnKF) and variational
4D-Var, address the same broad problem of hidden state estimation from observations over
time. Three structural differences distinguish the Yarncrawler framework. First, data
assimilation works in $\mathcal{X}$ rather than $\Xmod$, treating gauge redundancies as
noise rather than structure. Second, the regularizer in data assimilation is a prior
distribution expressing preference, not an admissibility condition expressing exclusion.
Third, data assimilation has no sheaf structure: there is no notion of local charts,
restriction maps, or cohomological obstruction, and no counterpart to the Identifiability
Theorem. In the linear Gaussian regime with trivial symmetry group, $\mathcal{C}$ reduces
to the Kalman update, making EnKF a degenerate special case.

\subsection{Inverse Rendering and Neural Radiance Fields}

Neural radiance fields (NeRF) and related inverse rendering methods reconstruct 3D scenes
from 2D images. The RGB projection operator $\Pi_{\mathrm{RGB}}$ of
Section~\ref{sec:projections} is the NeRF volume rendering integral. Three differences
are structural. NeRF has no dynamics: it is a static reconstruction with no dynamical
admissibility set $\Gamma$. NeRF has no RSVP field structure: the scene representation
is a learned radiance function, not a physical field triple. NeRF has no regularizer
encoding physical admissibility. The Yarncrawler framework strictly generalizes NeRF: RGB
is one chart of the cover, and NeRF is the degenerate case where only that chart is
available and dynamics are absent.

\subsection{Digital Twin Systems}

Digital twin literature proposes maintaining synchronized virtual replicas of physical
systems. The standard approach is object-ontological: the twin represents a collection of
objects whose properties are updated from sensor data. Yarncrawler departs in three ways.
The primary ontology is trajectory-based: objects emerge as persistent patterns in field
trajectories, not as primary entities. The twin is not updated by recognized objects but
by the consistency operator acting on the full quotient configuration space. Most
importantly, digital twins have no identifiability theory: there is no formal account of
when a twin uniquely determines its physical counterpart. The Yarncrawler Identifiability
Theorem provides exactly this.

\subsection{Categorical and Sheaf-Theoretic Semantics}

Sheaf theory has been applied to contextuality in quantum mechanics, to distributed
knowledge in multi-agent systems, and to compositional natural language semantics. The
Yarncrawler framework draws on these applications but departs in one crucial direction:
the sheaf here is over a physical configuration space with RSVP field structure, not over
an abstract semantic space. The stalks are Sobolev-class field configurations satisfying
physical equations. The repair morphisms are steps of a variational optimization. The
entropy identified with $H^1$ is a physical field, not an information-theoretic measure.
This grounding allows the framework to make contact with sensor physics and deployment
engineering while retaining the full generality of the categorical structure.

%% ============================================================
\section{Discussion: Closure as Epistemology}
\label{sec:epistemology}
%% ============================================================

\subsection{What Closure Means}

A closed Yarncrawler state is not a probabilistic estimate. It is not the most likely
world-state given the observations. It is the unique world-state that is structurally
demanded by the observations under the model class $\mathcal{A}$. The distinction has
operational consequences: a probabilistic estimate can always be revised by new
observations, while a closed state cannot be revised without either adding new sensor
modalities (killing new $H^1$ classes) or revising the model class. The closed state is,
in a precise sense, as certain as the model class allows.

\subsection{Constraint Primacy}

The framework operationalizes a structural commitment: constraints are prior to
completions. The world-state is not recovered by filling in a model template with
observations; it is the unique configuration surviving after all constraints have been
applied. The Sheaf-Variational Equivalence makes this precise: a world-state has meaning
(is uniquely identified) exactly when the sheaf induced by its projections has a unique
global section. This is constraint primacy as a mathematical theorem rather than a
philosophical position.

\subsection{The RSVP Fields as Constraint Carriers}

The RSVP field triple $(\Phi, \mathbf{v}, S)$ is not merely a convenient
parametrization of the configuration space. It is the structure that makes the regularizer
$\mathcal{R}$ non-arbitrary: the admissible set $\mathcal{A}$ is defined by the RSVP
field equations, which encode physical conservation laws, ecological succession
principles, and civilizational feasibility constraints. Without the RSVP structure, the
regularizer would be an arbitrary penalty and the closed state would be an artifact of
the choice of penalty. With it, the closed state is the unique configuration consistent
with the governing laws of the domain.

\subsection{Stigmergic Accumulation and the Refinement Loop}

Each step of the refinement loop of Section~\ref{sec:inversion} is a repair morphism in
the sheaf $\mathcal{S}_\Pi$: it reduces a projection residual (trivializes a cocycle) and
updates the field configuration toward closure. The accumulation of constraints over
time --- as more sensor observations arrive --- is the variational analogue of stigmergic
reinforcement: each new observation narrows the feasible set. This is the physical
realization of the self-refactoring dynamics that characterize the Yarncrawler framework
in its forward direction. The refinement loop converges when no more constraints can be
satisfied --- when the global section is unique and the yarn has been fully crawled.

\subsection{Limits of the Framework}

Three structural limits warrant acknowledgment.

\emph{Model class dependence}: closure is unique within $\mathcal{A}$, not in $\mathcal{X}$.
If $\mathcal{A}$ is misspecified --- if the true world-state lies outside the admissible
set --- the consistency operator converges to the closest admissible state, not the true
one. Non-existence of a global section is the falsification signal requiring model
revision.

\emph{Trajectory learning quality}: dynamical disambiguation depends on the quality of
$F(\cdot\,;\hat{\theta})$. If the operator class library is missing a dominant process
type, trajectory separation fails along the missing direction, detectable as residual
$|\mathcal{F}| > 1$ after full sensor coverage.

\emph{Computational tractability}: $\mathcal{C}$ is defined as a global minimizer of a
non-convex functional. The refinement loop approximates with local gradient steps.
Convergence to the global minimizer is guaranteed only up to initialization quality;
multiple encoder samples may be required.

\subsection{Forward and Inverse as Dual Inference}

The forward direction maintains coherence by repairing tears as they arise: it is the
epistemic activity of a system embedded in a world it partially controls. The inverse
direction reconstructs state from observations: it is the epistemic activity of a system
observing a world it does not control. The Yarncrawler framework is the unified theory of
both activities. The sheaf $\mathcal{S}_\Pi$ is the common object; the two directions are
two ways of relating to it. A fully realized system alternates between them:
reconstructing from observations, projecting forward, detecting model failure at
divergence, and repairing when failure is detected. This is the complete epistemic loop.

%% ============================================================
\section{Conclusion}
\label{sec:conclusion}
%% ============================================================

The Yarncrawler Identifiability Theorem and the Dynamical Disambiguation Theorem
characterize the conditions under which multi-modal sensor observations uniquely determine
a physical world-state. The Sheaf-Variational Equivalence shows that this characterization
is not merely a technical result about inverse problems: it is the variational realization
of a categorical structure in which sensor modalities are charts, observations are local
sections, and closure is the existence of a unique global section. The RSVP field triple
is the substrate in which these structures are physically grounded. The Morse structure of
the consistency functional, the Wasserstein gradient flow over world-state measures, and
the RSVP field evolution under constraint forcing are not three descriptions of the
framework --- they are one dynamical system in three coordinate systems.

A world-state is uniquely recoverable from multi-modal observations if and only if the
induced map from the quotient configuration space to the product of observation spaces is
injective on the feasible set, or if the dynamics of the world separate configurations
that static sensors cannot. In either case, closure corresponds simultaneously to
cohomology trivialization, entropy collapse, Fisher volume collapse, and convergence of the
Wasserstein gradient flow to a Dirac measure. These are not independent conditions; they
are one condition in four languages.

Three open problems remain. First, the extension of the identifiability theorem to the
stochastic case without the low-noise limit: does a version hold when $\sigma^2$ is not
small? Second, the learning theory of $F(\cdot\,;\hat{\theta})$: what sample complexity
is needed for the operator class library to ensure trajectory separation with high
probability? Third, the connection between the Fisher information geometry of the feasible
set and the computational complexity of the refinement loop: can convergence rate be
bounded in terms of the eigenspectrum of the Fisher metric at $[X^*]$?

The Yarncrawler system is not a sensing architecture. It is a theory of when the world
becomes legible --- when the observations accumulated by a sensor network, constrained by
physical law and dynamical history, leave no alternative configuration standing. That
condition is not an approximation target. It is a threshold: on one side, multiple
histories remain consistent with the evidence; on the other, only one does. The framework
identifies that threshold precisely and provides the operator that reaches it.

%% ============================================================
%% APPENDICES
%% ============================================================

\appendix

%% ============================================================
\section{Fréchet Manifold Structure of $\mathcal{X}$}
\label{app:frechet}
%% ============================================================

\subsection{Fréchet Manifolds}

A Fréchet manifold is a topological space locally homeomorphic to a Fréchet space (a
complete metrizable locally convex topological vector space) with smooth transition maps.
The configuration space $\mathcal{X} = H^1(\Omega) \times H^1(\Omega;\mathbb{R}^3) \times
L^2(\Omega)$ is a Hilbert space, hence also a Banach space and a Fréchet space, with its
product inner product. The Fréchet framing is introduced because extensions to more
general field configurations (distributional fields, fields on non-compact domains) may
require genuine Fréchet rather than Hilbert structure.

\subsection{The Quotient Topology}

The symmetry group $G$ acts on $\mathcal{X}$ smoothly, properly, and isometrically.
Under these conditions the quotient $\mathcal{X}/G$ is Hausdorff. The metric on $\Xmod$
is induced from $\mathcal{X}$ via $d([X],[X']) = \inf_{g \in G} \|X - gX'\|_{\mathcal{X}}$.

The observational quotient $\Xmod$ (by $\sim$) is in general coarser than $\mathcal{X}/G$.
They coincide when the projection family is complete (Proposition~\ref{prop:completeness}),
because then every sensor-invisible direction is also $G$-equivalent.

\subsection{Tangent Bundle of $\Xmod$}

The tangent bundle $T\Xmod$ is the quotient of $T\mathcal{X}$ by the action of $TG$. When
$G$ acts freely, $T\Xmod \cong T\mathcal{X}/TG$ is a vector bundle. The admissible tangent
cone $\mathcal{T}_{[X]}\Gamma^0$ (Definition~6.5 in the main text) is a linear subspace
of $T_{[X]}\Xmod$ of dimension at most $p = \dim\Theta$ (Proposition~6.6).

\subsection{Geodesic Convexity on $\Xmod$}

Geodesics on $\Xmod$ are the images of straight lines in $\mathcal{X}$ under the
quotient map. The admissible set $\mathcal{A}$ is geodesically convex when its preimage
in $\mathcal{X}$ is convex, which holds by definition of $\mathcal{A}$ as the
sublevel set $\{\mathcal{R}(X) < \infty\}$ of the convex functional $\mathcal{R}$.
This justifies Assumption~\ref{ass:admissibility}.

%% ============================================================
\section{Full Proof of the Identifiability Theorem}
\label{app:identifiability}
%% ============================================================

\subsection{Functional-Analytic Setting}

The proof of Theorem~\ref{thm:identifiability} relies on existence and uniqueness of
minimizers of the consistency functional on $\Xmod$. We introduce additional assumptions
for the infinite-dimensional setting.

\begin{assumption}[Banach Realization, restated]
There exists a Banach space $\mathcal{B}$ and continuous embedding
$\iota : \mathcal{X} \hookrightarrow \mathcal{B}$ such that $\mathcal{A}$ is closed and
convex in $\mathcal{B}$, each $\Pi_i$ extends to a bounded linear operator on $\mathcal{B}$,
and $\mathcal{R}$ extends to a proper, lower semicontinuous functional on $\mathcal{B}$.
In practice, $\mathcal{B} = \mathcal{X}$ itself.
\end{assumption}

\begin{assumption}[Geodesic Convexity, restated]
$\mathcal{A}$ is geodesically convex with respect to the metric induced from $\mathcal{B}$,
and $\mathcal{R}$ is strictly convex along geodesics in $\mathcal{A}$:
\[
\mathcal{R}(\gamma(t)) < (1-t)\mathcal{R}([X]) + t\mathcal{R}([X'])
\quad \forall [X] \ne [X'] \in \mathcal{A},\; t \in (0,1).
\]
\end{assumption}

\begin{assumption}[Coercivity]
$f([X]) = \sum_i \|\Pi_i(X)-y_i\|^2 + \mathcal{R}(X)$ satisfies
$\|X_n\|_{\mathcal{B}} \to \infty \Rightarrow f([X_n]) \to \infty$.
\end{assumption}

\subsection{Existence of a Minimizer}

\begin{proposition}
Under the standing assumptions, $f$ attains a minimum on $\mathcal{A}$.
\end{proposition}
\begin{proof}
Let $\{[X_n]\}$ be a minimizing sequence; coercivity gives boundedness in $\mathcal{B}$.
Reflexivity of $\mathcal{B}$ (if Hilbert) yields a weakly convergent subsequence
$X_{n_k} \rightharpoonup X_*$. Closedness of $\mathcal{A}$ gives $[X_*] \in \mathcal{A}$.
Weak continuity of $\Pi_i$ (bounded linear) and lower semicontinuity of $\mathcal{R}$
give $f([X_*]) \le \liminf f([X_{n_k}])$. Hence $[X_*]$ is a minimizer. $\square$
\end{proof}

\subsection{Uniqueness of the Minimizer}

\begin{proposition}
Under the standing assumptions, the minimizer is unique.
\end{proposition}
\begin{proof}
Suppose $[X],[X']$ are distinct minimizers. By geodesic convexity, there exists a
geodesic $\gamma$ in $\mathcal{A}$ connecting them. The projection residual term is
convex along $\gamma$ (by linearity of $\Pi_i$). The regularizer is strictly convex along
$\gamma$ by assumption. Hence $f(\gamma(t)) < (1-t)f([X]) + tf([X']) = \inf f$ for
$t \in (0,1)$, contradicting minimality. $\square$
\end{proof}

\subsection{Passage to the Quotient}

\begin{proposition}
If $[X] = [X']$ in $\Xmod$, then $f([X]) = f([X'])$. Conversely, distinct equivalence
classes yield distinct objective values whenever $\Ftil$ is injective.
\end{proposition}
\begin{proof}
Observational equivalence means $\Pi_i(X) = \Pi_i(X')$ for all $i$, so residual terms
are equal; $\mathcal{R}$ is defined on equivalence classes, so regularizer terms are
equal. Injectivity of $\Ftil$ means distinct classes differ in at least one projection,
so residual terms differ. $\square$
\end{proof}

\subsection{Full Proof of Theorem~\ref{thm:identifiability}}

\begin{proof}[Full proof]
\textit{Well-definedness}: by the existence and uniqueness propositions, $\mathcal{C}$ is
single-valued on $\mathcal{A}$.

\textit{Fixed points = minimizers}: $[X^*] = \mathcal{C}([X^*])$ iff $[X^*]$ minimizes
$f$ over $\mathcal{A}$.

\textit{Injectivity implies uniqueness}: if $\Ftil$ is injective on $\mathcal{F}$,
$\mathcal{F}$ contains at most one class. By non-emptiness, $[X^*] \in \mathcal{F}$, so
$|\mathcal{F}| = 1$. At $[X^*]$, all residuals are within tolerance; strict convexity
prevents alternative minimizers. Hence $[X^*]$ is the unique fixed point.

\textit{Non-injectivity implies non-uniqueness}: if $\Ftil$ is not injective on
$\mathcal{F}$, distinct $[X],[X'] \in \mathcal{F}$ have identical projections. Both are
candidates for the minimizer. If $\mathcal{R}([X]) = \mathcal{R}([X'])$, strict
convexity fails (contradiction). If $\mathcal{R}([X]) \ne \mathcal{R}([X'])$, the
larger-regularizer point is not a minimizer---but then it is not in $\mathcal{F}$ as
defined, contradicting the assumption. Either way we reach a contradiction with
uniqueness. $\square$
\end{proof}

%% ============================================================
\section{Full Proof of the Dynamical Disambiguation Theorem}
\label{app:disambiguation}
%% ============================================================

\subsection{Measure-Theoretic Framework}

\begin{assumption}[Polish Structure]
$\mathcal{A}$ with the metric induced from $\mathcal{B}$ is a Polish space (complete,
separable metric).
\end{assumption}

\begin{definition}[Markov Kernel, precise]
$T_\Delta : \mathcal{A} \to \mathcal{P}(\mathcal{A})$ is a Markov kernel: for each
$[X]$, $T_\Delta([X],\cdot)$ is a Borel probability measure on $\mathcal{A}$; for each
Borel $B \subseteq \mathcal{A}$, $[X] \mapsto T_\Delta([X],B)$ is measurable.
\end{definition}

\subsection{Convergence to the Deterministic Generator}

\begin{assumption}[Generator Convergence, precise]
For all bounded continuous $\varphi : \mathcal{A} \to \mathbb{R}$:
\[
\int \varphi(X')\,T_\Delta([X],dX') = \varphi([X] + \Delta F([X];\theta)) + O(\Delta^{3/2}).
\]
\end{assumption}

This is weak convergence $T_\Delta([X],\cdot) \Rightarrow \delta_{[X]+\Delta F([X];\theta)}$
at rate $\Delta^{3/2}$, corresponding to a diffusion with noise variance $O(\Delta)$.

\subsection{Stability of Trajectory Separation Under Noise}

\begin{definition}[Uniform Separation Margin]
A pair $[X],[X'] \in \Omega_{\mathrm{obs}}$ is \emph{uniformly trajectory-separated} if
there exist $\tau > 0$ and $\delta > 0$ such that
\[
\left\| \Ftil(\Phi^\tau_F([X])) - \Ftil(\Phi^\tau_F([X'])) \right\| \ge \epsilon + \delta.
\]
\end{definition}

\begin{lemma}
Under generator convergence, uniformly trajectory-separated pairs remain distinguishable
under stochastic evolution with probability $\ge 1 - C\Delta/\delta^2$.
\end{lemma}
\begin{proof}
By Chebyshev applied to the $O(\Delta)$ noise variance, stochastic trajectories deviate
from deterministic by $O(\sqrt{\Delta})$ with high probability. For deviation $\delta/2$:
$\mathbb{P}(\|X_\tau^\Delta - X_\tau\| > \delta/2) \le C\Delta/\delta^2$.
With margin $\delta$, the two stochastic trajectories remain separated by $\ge \epsilon$
with probability $\ge 1 - 2C\Delta/\delta^2$. $\square$
\end{proof}

\subsection{Full Proof of Theorem~\ref{thm:disambiguation}}

\begin{proof}[Full proof]
Let $[X],[X'] \in \Omega_{\mathrm{obs}} \cap \Gamma^0$ be distinct and uniformly
trajectory-separated at time $\tau$.

\textit{Step 1}: Both lie on integral curves of $F(\cdot\,;\hat{\theta})$ by definition of
$\Gamma^0$.

\textit{Step 2}: At time $\tau$, their projections differ by $\ge \epsilon + \delta > \epsilon$.

\textit{Step 3}: By the noise lemma, stochastic evolutions remain separated with high
probability for small $\Delta$.

\textit{Step 4}: At time $\tau$, the true observation $y(\tau)$ has
$\|y_i(\tau) - \Pi_i(X^*(\tau))\| \le \epsilon_i$. At most one of the two separated
candidates lies within $\epsilon_i$ of $y_i(\tau)$.

\textit{Step 5}: The true state $[X^*(\tau)] \in \Omega_{\mathrm{obs}}(y(\tau))$ by
non-emptiness applied at time $\tau$. Hence the false candidate is eliminated.

\textit{Step 6}: Applying to all pairs eliminates all but $[X^*]$. Hence
$|\mathcal{F}| = 1$. $\square$
\end{proof}

%% ============================================================
\section{Fisher Information Geometry of the Feasible Set}
\label{app:fisher}
%% ============================================================

\subsection{Purpose and Position}

Section~\ref{sec:closure} states that closure corresponds to vanishing information volume
of the feasible set. This appendix provides the formal construction underlying that claim.
The central point is that the projection family $\{\Pi_i\}$, together with noise models
$\{\Sigma_i\}$, induces a Riemannian metric on the admissible quotient space. The metric
measures distinguishability of nearby world-states under the available sensor suite.
Closure occurs when all remaining admissible directions are eliminated.

\subsection{Observation Likelihood on the Quotient Space}

For each modality $i$, assume noise $\eta_i \sim \mathcal{N}(0, \Sigma_i)$ so the
conditional density is
\[
p_i(y_i \mid X) \propto \exp\!\left(
-\tfrac{1}{2}\|y_i - \Pi_i(X)\|_{\Sigma_i^{-1}}^2
\right).
\]
Assuming conditional independence across modalities, the full observation log-likelihood is
\[
-\ell(X; y) = \tfrac{1}{2}\sum_{i \in \mathcal{I}} \|y_i - \Pi_i(X)\|_{\Sigma_i^{-1}}^2.
\]
This is well-defined on $\mathcal{A} = \Xmod$ because each $\Pi_i$ is constant on
$\sim$-equivalence classes. The negative log-likelihood is precisely the residual sum
in the consistency operator objective.

\subsection{The Fisher Information Metric}

\begin{definition}[Fisher Information Metric]
The Fisher information metric $g_F$ on $\mathcal{A}$ is the symmetric bilinear form
\[
g_F(X)(u, v) = \sum_{i \in \mathcal{I}}
\left\langle D\Pi_i(X)u,\; \Sigma_i^{-1} D\Pi_i(X)v \right\rangle_{\mathcal{Y}_i},
\]
where $D\Pi_i(X) : T_X\mathcal{A} \to \mathcal{Y}_i$ is the derivative of $\Pi_i$ at $[X]$.
\end{definition}

This is a positive semidefinite tensor. It is positive definite exactly when the local
model is identifiable: no nonzero tangent direction leaves the induced observation law
unchanged to first order. The null space of $g_F$ is therefore the infinitesimal version
of projection degeneracy:
\[
g_F(X)(v,v) = 0 \;\Longleftrightarrow\; D\Pi_i(X)v = 0 \;\forall i.
\]
The negative log-likelihood Hessian equals $g_F$ in the Gaussian case, linking the
Fisher metric directly to the curvature of the consistency operator objective.

\subsection{Metric Monotonicity Under Sensor Addition}

\begin{proposition}[Metric Monotonicity]
\label{prop:metric_monotone}
Adding a new sensor modality $j$ yields
$g_F^{\mathrm{new}}(X) \succeq g_F(X)$ at every $[X] \in \mathcal{A}$.
\end{proposition}
\begin{proof}
$g_F^{\mathrm{new}}(X)(u,u) = g_F(X)(u,u) +
\|D\Pi_j(X)u\|_{\Sigma_j^{-1}}^2 \ge g_F(X)(u,u)$.
$\square$
\end{proof}

This is the infinitesimal counterpart of Proposition~\ref{prop:contraction}: new sensors
cannot reduce local distinguishability. They fail to improve it only when their derivative
contribution lies entirely in the directions already constrained.

\subsection{Restriction to the Feasible Set}

\begin{definition}[Restricted Fisher Metric]
The restricted Fisher metric on $\mathcal{F}$ is
\[
g_F^{\mathcal{F}}(X) = g_F(X)\big|_{T_X\mathcal{F} \times T_X\mathcal{F}}.
\]
\end{definition}

When $\mathcal{F}$ is locally a smooth embedded submanifold of $\mathcal{A}$ (which holds
on regular strata under generic regularity assumptions), $g_F^{\mathcal{F}}$ defines a
Riemannian metric on those strata.

\begin{definition}[Information Volume]
\[
\mathrm{Vol}_F(\mathcal{F}) = \int_{\mathcal{F}_{\mathrm{reg}}} d\mathrm{Vol}_{g_F^{\mathcal{F}}}.
\]
\end{definition}

Large information volume indicates many locally distinguishable admissible world-states.
Zero volume indicates closure, subject to the nondegeneracy condition below.

\subsection{Dynamical Lifting of Null Directions}

Define the time-extended projection $\Pi_i^{(\tau)}(X) = \Pi_i(\Phi_F^\tau(X))$ and its
induced metric
\[
g_F^{(\tau)}(X)(u,v)
= \sum_i \left\langle D(\Pi_i \circ \Phi_F^\tau)(X)u,\;
\Sigma_i^{-1} D(\Pi_i \circ \Phi_F^\tau)(X)v \right\rangle.
\]
By the chain rule, $D(\Pi_i \circ \Phi_F^\tau)(X) = D\Pi_i(\Phi_F^\tau(X)) \circ D\Phi_F^\tau(X)$.

\begin{proposition}[Metric Form of Dynamical Disambiguation]
\label{prop:dynamical_lifting}
Suppose $g_F(X)(v,v) = 0$ but $g_F^{(\tau)}(X)(v,v) > 0$ for some $\tau > 0$.
Then $v$ is a static null direction that is lifted into the observable range by the
dynamics $\Phi_F^\tau$. This is the differential-geometric form of
Theorem~\ref{thm:disambiguation}: two configurations indistinguishable at $t = 0$
become distinguishable at $t = \tau$.
\end{proposition}

Time is therefore not an auxiliary feature---it is the mechanism by which null directions
of the static Fisher metric are eliminated.

\subsection{Closure as Vanishing Information Volume}

\begin{theorem}[Closure Implies Vanishing Information Volume]
\label{thm:closure_volume}
If $\mathcal{F} = \{[X^*]\}$, then $\mathrm{Vol}_F(\mathcal{F}) = 0$.
\end{theorem}
\begin{proof}
A singleton has zero-dimensional regular part; the Riemannian volume form integrates to
zero. $\square$
\end{proof}

\begin{proposition}
\label{prop:volume_converse}
$\mathrm{Vol}_F(\mathcal{F}) = 0$ with $g_F^{\mathcal{F}}$ nondegenerate on the regular
part of $\mathcal{F}$ implies $\dim_{\mathrm{eff}}(\mathcal{F}) = 0$.
\end{proposition}
\begin{proof}
Under nondegeneracy, positive dimension implies positive volume. $\square$
\end{proof}

\begin{corollary}[Closure Criterion via Information Geometry]
\label{cor:closure_criterion}
If $g_F^{\mathcal{F}}$ is nondegenerate on the regular part of $\mathcal{F}$, then
\[
\mathcal{F} = \{[X^*]\}
\;\Longleftrightarrow\;
\dim_{\mathrm{eff}}(\mathcal{F}) = 0
\;\Longleftrightarrow\;
\mathrm{Vol}_F(\mathcal{F}) = 0.
\]
\end{corollary}

The nondegeneracy assumption excludes the case where $\mathcal{F}$ has positive dimension
but the metric degenerates on it (unresolved indistinguishability rather than genuine
uniqueness). This is the exact condition Section~\ref{sec:closure} relies on.

\subsection{Reconstruction Efficiency}

Let $\mathcal{F}_n$ denote the feasible set after $n$ sensor modalities or constraint
updates. Define:

\begin{definition}[Reconstruction Efficiency]
\[
\mathcal{E}_n = 1 - \frac{\mathrm{Vol}_F(\mathcal{F}_{n+1})}{\mathrm{Vol}_F(\mathcal{F}_n)}
\in [0,1].
\]
\end{definition}

Values near $0$ indicate the new sensor or update contributes little additional
discrimination. Values near $1$ indicate near-total collapse of residual admissible
uncertainty. A continuous-time analogue is
$\mathcal{E}(t) = -\frac{d}{dt}\log \mathrm{Vol}_F(\mathcal{F}(t))$,
which behaves as an entropy-contraction rate and is often more stable than raw volume
differences.

The optimal sensor placement criterion follows: the next sensor $j$ should be placed
where the directional derivative of $\mathcal{E}$ is maximized, which is determined by
the eigenspectrum of $g_F$ in the remaining unclosed directions.

\subsection{The Effective Metric Near Closure}

Near a nondegenerate minimizer $[X^*]$, the Hessian of the consistency objective
$J(X) = \sum_i \|y_i - \Pi_i(X)\|_{\Sigma_i^{-1}}^2 + \mathcal{R}(X)$ is
\[
\nabla^2 J(X^*) = g_F(X^*) + \nabla^2 \mathcal{R}(X^*).
\]
This effective metric combines statistical distinguishability (Fisher geometry) with
admissibility curvature (regularizer). The convergence rate of the refinement loop of
Section~\ref{sec:inversion} is governed by the condition number
$\kappa = \lambda_{\max}/\lambda_{\min}$ of $\nabla^2 J(X^*)$. Directions with small
Fisher information and small regularizer curvature correspond to slow convergence; dense
multi-modal sensor networks increase $\lambda_{\min}$ and hence accelerate convergence
to closure.

\subsection{Entropy Identification}

The identification stated in Section~\ref{sec:equivalence},
\[
\log \mathrm{Vol}_F(\mathcal{F})
\;\longleftrightarrow\;
\log \dim H^0(\mathcal{U}, \mathcal{S}_\Pi)
\;\longleftrightarrow\;
S_{\mathrm{RSVP}}(X^*),
\]
is now fully grounded. The first identification holds in the Gaussian regime because
$\dim H^0$ counts linearly independent global sections, which scales with the information
volume of the feasible set. The second uses the RSVP entropy field equation: at the fixed
point $[X^*]$, $S_{\mathrm{RSVP}}$ is the integral of local entropy production over
$\Omega$, which under the consistency operator objective equals the log-volume of
remaining uncertainty in the observation-consistent set. The RSVP entropy field is the
local density of cohomological obstruction; its minimization under $\mathcal{C}$ is the
physical expression of sheaf repair.

%% ============================================================
\section{Operator Class Library}
\label{app:operators}
%% ============================================================

\subsection{Structure of the Library}

The library consists of $K$ named process generators $\{F_k\}_{k=1}^K$, each a learned
vector field on $\mathcal{A}$ parametrized by $\theta_k \in \Theta_k$. The full dynamics
decompose as
\[
F([X];\theta) = \sum_{k=1}^K \alpha_k([X])\, F_k([X];\theta_k),
\]
with activation weights $\alpha_k([X]) \ge 0$ satisfying $\sum_k \alpha_k = 1$. Each
$F_k$ is learned from before/after pairs in which process $k$ is dominant, identified by
its characteristic sensor signature.

\subsection{Ecological Process Operators}

\emph{Forest succession} ($F_{\mathrm{succ}}$): slow increase in LiDAR canopy height,
gradual RGB reflectance shift from grass-like to canopy-like, decrease in audio spectral
entropy as bird diversity changes with succession stage. Time scale: years to decades.

\emph{Disturbance and recovery} ($F_{\mathrm{dist}}$): rapid decrease in LiDAR height
and RGB NDVI followed by slow recovery; audio signature is absence of characteristic
species. Time scale: event-driven (days), recovery (years).

\emph{Hydrological dynamics} ($F_{\mathrm{hydro}}$): correlated changes in infrared
(soil moisture affects surface temperature), LiDAR (inundation changes geometry), and
audio (water flow sounds). Time scale: seasonal to event-driven.

\emph{Phenological cycling} ($F_{\mathrm{pheno}}$): regular seasonal oscillations in RGB
(green-up/senescence), infrared (latent heat flux), and audio (bird song seasonality).
Time scale: annual.

\subsection{Civilizational Process Operators}

\emph{Urban expansion} ($F_{\mathrm{urban}}$): progressive increase in LiDAR building
density and height, RGB shift to impervious surface reflectance, audio increase (traffic,
machinery). Time scale: years to decades.

\emph{Infrastructure decay} ($F_{\mathrm{decay}}$): gradual LiDAR surface roughness
change, RGB spectral shift, audio entropy increase (structural noise). Time scale: years.

\emph{Agricultural cycling} ($F_{\mathrm{agri}}$): regular RGB oscillations (crop growth
stages), infrared ET changes, LiDAR canopy height variation, audio machinery activity.
Time scale: seasonal.

\emph{Demographic change} ($F_{\mathrm{demo}}$): correlated changes in audio activity
levels, RGB night-light proxies, and infrastructure use patterns. Time scale: years.

\subsection{Learning Protocol}

Each $F_k$ is learned from before/after pairs as follows. First, identify pairs where
process $k$ is dominant by classifying the observation delta against the existing library.
Second, fit $\theta_k$ by minimizing the forward-projection discrepancy:
\[
\hat{\theta}_k = \arg\min_{\theta_k} \sum_{n \in \mathcal{D}_k} \sum_i
\left\|\Pi_i(\Phi^\tau_{F_k}([X_b^{(n)}];\theta_k)) - y_a^{(n)}\right\|^2.
\]
Third, validate on held-out pairs by checking that
$\Theta_{(y_b,y_a,\tau)} \cap \Theta_k \ne \emptyset$ for pairs labeled as process $k$.

\subsection{Library Extensibility}

The library is open: new process types $F_{K+1}$ are added by collecting before/after
pairs exhibiting the new process, fitting the generator as above, and appending to the
library. The identifiability and disambiguation theorems are unaffected by library
extension: they hold for any $F$ with Lipschitz generator. The effect of extension is to
reduce $\Gamma^0_{\mathrm{indist}}$ --- the set of configurations not separated by
trajectory learning --- which by Corollary~\ref{cor:joint_identifiability} increases the
rate of closure.

%% ============================================================
\section{Implementation File Structure}
\label{app:implementation}
%% ============================================================

This appendix documents the reference implementation of the Yarncrawler consistency
operator as a stratified filesystem. The architecture maps directly onto the theoretical
structure: state definition corresponds to the configuration space $\mathcal{X}$,
morphism algebra corresponds to the operator class library, controller evolution
corresponds to iterative application of $\mathcal{C}$, and quotient construction
corresponds to the identification of observationally equivalent world-states.

\subsection{Structural Overview}

The implementation is organized into seven layers, each corresponding to a structural
role in the framework.

\begin{center}
\begin{tabular}{ll}
\toprule
\textbf{Layer} & \textbf{Theoretical Role} \\
\midrule
State layer & Configuration space $\mathcal{X}$, field types, domains \\
Morphism layer & Operator class library $\{F_k\}$, composition algebra \\
Transfer layer & Domain transfer, functoriality checking \\
Controller layer & Iterative application of $\mathcal{C}$, decision policy \\
Transformation layer & Natural transformations between controllers \\
Fixed-point layer & Closure detection, stability tracking \\
Quotient layer & Equivalence classes of world-states \\
\bottomrule
\end{tabular}
\end{center}

\subsection{Core Infrastructure}

Shared utilities used by all layers:
\begin{verbatim}
common/
  common_io.py          # I/O primitives, serialization, logging
\end{verbatim}

\subsection{State Layer}

Defines the plenum field configuration and admissible types.
\begin{verbatim}
state/
  initial_state.json    # Initial field configuration (Phi, v, S)
  topology.json         # Domain geometry, boundary conditions
  typing_rules.yaml     # Type constraints on field components

  domain_A.json         # Source domain specification
  domain_B.json         # Target domain specification

  transfer_A_to_B.json  # Cross-domain transfer morphism
\end{verbatim}

\subsection{Log Layer}

Accumulates the trajectory history that the disambiguation theorem requires.
\begin{verbatim}
logs/
  history.jsonl           # Full state history (before/after pairs)
  controller_trace.jsonl  # Controller decision trace
  class_transforms.jsonl  # Equivalence class transitions
\end{verbatim}

\subsection{Morphism Layer}

Implements the operator class library: storage, generation, composition, transfer.
\begin{verbatim}
morphisms/
  registry/
    index.json            # Global operator index
    entries/
      morph_*.json        # Individual operator specifications

  transferred/
    *.json                # Operators transferred across domains

python/
  morphism_apply.py       # Apply operator to field configuration
  compose_morphisms.py    # Categorical composition of operators
  discover_compositions.py # Find composable operator pairs
  abstract_morphism.py    # Abstract operator from trajectory data

  admit_morphism.py       # Admit new operator to registry
  update_registry_stats.py # Track operator usage statistics
  prune_registry.py       # Remove underused operators
  query_registry.py       # Retrieve operators by type/signature

  transfer_morphism.py    # Transfer operator across domains
  transfer_registry_slice.py # Transfer operator subset
  check_functoriality.py  # Verify naturality conditions

  check_composable.py     # Type-check operator composition
  type_path_resolver.py   # Resolve type paths in composition
  propose_type_bridge.py  # Propose bridging morphisms

  score_rollout.py        # Score operator on trajectory rollout
\end{verbatim}

\subsection{Controller Layer}

Implements the iterative consistency operator $\mathcal{C}$ and decision policies.
\begin{verbatim}
controllers/
  registry/
    beam.json       # Beam search controller
    dp.json         # Dynamic programming controller
    learned.json    # Learned neural controller
    spatial.json    # Spatially-indexed controller
    repair.json     # Sheaf repair controller

  state/
    controller_state.json  # Current controller parameters

  transformations/
    beam_to_learned.json   # Natural transformation between controllers

  quotient/
    classes.json            # Equivalence classes of controllers
    class_0001.json         # Class representative and members
    class_0002.json
    representatives.json    # Selected representatives per class

python/
  apply_controller_transform.py   # Apply natural transformation
  learn_controller_transform.py   # Learn transformation from data

  iterate_controller.py     # One step of consistency operator
  controller_distance.py    # Distance between controller states
  check_fixed_point.py      # Test for closure (fixed point)
  update_stability.py       # Track convergence to closure

  controller_equivalence.py # Identify equivalent controllers
  build_controller_classes.py # Construct equivalence classes
  select_representatives.py  # Choose class representatives

  build_controller_quotient.py  # Construct quotient structure
  split_classes.py              # Refine equivalence classes
  select_controller_class.py    # Choose active class
  log_class_transform.py        # Record class transitions
  canonicalize_controller.py    # Map to canonical representative

  meta_controller.py        # Orchestrate controller selection
\end{verbatim}

\subsection{Engine Layer}

Executable entry points that orchestrate the full refinement loop.
\begin{verbatim}
engine/
  apply_morphism.sh            # Apply operator to current state
  run_transfer.sh              # Execute domain transfer
  run_controller_fixed_point.sh # Run until closure detected
  run_class_controller.sh      # Run with class-level controller

run.sh                         # Top-level execution entry point
\end{verbatim}

\subsection{Correspondence to the Theoretical Framework}

The implementation instantiates the theoretical structure as follows. The
\texttt{state/} layer defines $\mathcal{X}$ and $\mathcal{A}$. The
\texttt{morphisms/} layer is the operator class library $\{F_k\}$ of
Appendix~\ref{app:operators}. The \texttt{logs/history.jsonl} file accumulates the
before/after trajectory pairs used to learn $F(\cdot;\hat{\theta})$. The
\texttt{iterate\_controller.py} script implements one step of $\mathcal{C}$.
The \texttt{check\_fixed\_point.py} script implements the closure detection criterion
of Section~\ref{sec:closure}: stability of $[X^*]$ under perturbation of $\epsilon_i$.
The \texttt{controllers/quotient/} directory implements the identification of
observationally equivalent world-states, corresponding to the passage to $\Xmod$.
The \texttt{check\_functoriality.py} script verifies that domain transfers respect
the naturality conditions of the functor $\Psi$ established in
Theorem~\ref{thm:equivalence}.

%% ============================================================
\section{The Constellation Construction as Iterated Constraint Collapse}
\label{app:filtration}
%% ============================================================

\subsection{The Initial Object}

Let $\mathcal{M}_0$ denote the space of all syntactically valid modules:
\[
\mathcal{M}_0 = \{ \texttt{module.yaml} \text{ satisfying schema} \}
\]
At this stage no global compatibility is enforced. The space is combinatorially large
and admits arbitrary configurations. Each subsequent stage reduces admissible degrees of
freedom by introducing new compatibility conditions. The resulting object is a tower of
constraint refinements converging to a unique global section.

\subsection{Stage-wise Refinement as Endofunctors}

Each stage $k$ defines a restriction operator:
\[
\mathcal{R}_k : \mathcal{M}_{k-1} \to \mathcal{M}_k, \qquad
\mathcal{M}_k = \{ X \in \mathcal{M}_{k-1} \mid X \text{ satisfies constraints of stage } k \}
\]
The full system is the intersection:
\[
\mathcal{M}_* = \bigcap_{k=1}^{13} \mathcal{M}_k
\]
The 13 stages are not a pipeline. They are a descending filtration:
\[
\mathcal{M}_0 \supset \mathcal{M}_1 \supset \cdots \supset \mathcal{M}_{13}
\]
and the development history documented in Appendix~\ref{app:implementation} through
Appendix~\ref{app:controller_quotient} is the construction of this filtration one level
at a time.

\subsection{Sheaf Enforcement (Stages 1--4)}

At Stage 1, the module space is lifted into a sheaf:
\[
\mathcal{S} : \mathcal{U}^{\mathrm{op}} \to \mathbf{Set}
\]
where modules are local sections and the constellation's global section corresponds to
the deployed system. The cover $\mathcal{U}$ is the stratum partition: each stratum
(Sensory, Intellect, Understanding, Decision, Reason, Will, Presentation, Thought)
defines one chart.

Stages 2--4 enforce the sheaf condition by eliminating non-gluing configurations:

Stage 2 detects cocycles: \texttt{check\_commutativity.sh} identifies pairs of modules
whose interface contracts are incompatible, emitting structured obstruction objects
corresponding to $c_{ij} \in H^1(\mathcal{U}, \mathcal{S})$.

Stage 3 expands the restriction maps by introducing a type DAG over field spaces
($H^1(\Omega) \hookrightarrow L^2(\Omega) \hookrightarrow \mathrm{Scalar}$),
replacing one-step compatibility checks with multi-hop path resolution.

Stage 4 constructs coboundaries: \texttt{propose\_bridge.sh} and
\texttt{apply\_repair.sh} insert bridge modules that trivialize detected cocycles.

After Stage 4, the system has driven $H^1(\mathcal{U}, \mathcal{S}) \to 0$ or has
produced a structured report identifying which obstructions are irreducible.

\subsection{Emergence of the Morphism Algebra (Stages 5--7)}

Stages 5--7 lift the system from object space to morphism space. Define:
\[
\mathcal{O} = \{ f : X \to X \mid f(\Phi, \mathbf{v}, S) = (a_\Phi\Phi + b_\Phi,\, a_v\mathbf{v} + \mathbf{b}_v,\, a_S S + b_S) \}
\]
the space of affine RSVP operators. Stage 5 introduces the morphism kernel and
normalizes all operators into this common form. Stage 6 introduces the candidate
generation loop: mutation, composition, evaluation, and the persistent registry.
Stage 7 moves the vector field $\mathbf{v}$ from node-based coordinates to edge-based
flux, making divergence a proper discrete conservation law and enabling
geometry-respecting operator semantics.

These stages impose closure under composition, compatibility with RSVP field structure,
and conservation laws on graph topology. The system transitions from a static constraint
set to a dynamical algebra: $\mathcal{M}_7$ is not just a set of valid configurations
but a set of valid configuration-evolution pairs.

\subsection{Variational Selection and Trajectory Space (Stages 8--10)}

Stages 8--10 introduce trajectory dependence. Define the rollout functional:
\[
\mathcal{J}(f; X_0) = \sum_{t=0}^T \mathcal{L}(f_t(X_t))
\]
where $X_{t+1} = F(f_t(X_t), X_t)$ is the coupled field dynamics and $\mathcal{L}$
integrates sensor residuals, divergence penalties, boundary leakage, and semantic
property verification. The admissible operator set becomes:
\[
\mathcal{O}_{\mathrm{adm}} = \{ f \mid \mathcal{J}(f) \text{ minimized under constraints} \}
\]

Stage 8 implements this functional as a gated rollout scorer. Stage 9 spatializes the
controller: instead of a single operator applying to the whole field, a weight field
$w : \Omega \to \Delta(\mathcal{O})$ assigns a distribution over operators to each node,
smoothed over the topology. Stage 10 introduces local morphism generation conditioned on
regional obstruction signatures: each region of the field can generate operators targeted
at its specific failure mode, and successful local repairs enter the global registry.

The effective dimension of $\mathcal{M}_{10}$ is significantly lower than
$\mathcal{M}_7$: not only must configurations satisfy static compatibility conditions,
they must also lie on admissible trajectories under the selected operators.

\subsection{Typed Category and Compositional Structure (Stage 11)}

Stage 11 introduces higher-order structure. The operator space becomes a typed category:
\[
\mathcal{O} : \mathcal{T} \to \mathcal{T}
\]
where objects are field types and morphisms are RSVP operators. Composition is
constrained by:
\[
f \circ g \;\text{defined iff}\; \mathrm{cod}(g) \leadsto \mathrm{dom}(f)
\]
via a path in the type DAG. The type path resolver finds chains of admissible embeddings,
and the composition engine (\texttt{python/compose\_morphisms.py}) verifies compatibility
before any chain is formed.

Composition discovery (\texttt{python/discover\_compositions.py}) identifies empirically
successful operator chains from the history log, and abstraction
(\texttt{python/abstract\_morphism.py}) compresses these into single operators, growing
the registry with reusable primitives. At this stage the operator space is no longer
flat: it has algebraic structure, and $\mathcal{M}_{11}$ is the intersection of
configuration space with this structured operator algebra.

\subsection{Functorial Transfer (Stage 12)}

Stage 12 introduces inter-domain structure. Let $\mathcal{D}_1, \mathcal{D}_2$ be
domain categories defined by their topology, sensor bundle, and field type system. A
transfer map defines a functor:
\[
\mathcal{F} : \mathcal{D}_1 \to \mathcal{D}_2, \qquad
\mathcal{F}(f \circ g) \approx \mathcal{F}(f) \circ \mathcal{F}(g)
\]
Deviation from functoriality defines a transfer error, scored by
\texttt{python/check\_functoriality.py}. Transferred operators enter the target registry
with a score discount proportional to that error.

This stage does not reduce the effective dimension of $\mathcal{M}_{12}$ within a single
domain. Its role is different: it introduces a relation of equivalence between admissible
configurations across domains, so that the system's learned structure can be reused
rather than rebuilt from scratch in each new environment.

\subsection{Controller Quotient (Stage 13)}

Stage 13 lifts the system to the highest level of abstraction. Controllers are functors:
\[
\mathcal{K} : \mathcal{X} \to \mathcal{O}
\]
mapping field states to operator sequences. Natural transformations $\eta : \mathcal{K}_1
\Rightarrow \mathcal{K}_2$ define equivalences between controllers. The quotient
$\mathcal{K}/{\sim}$ collapses redundant controllers into equivalence classes, and the
system operates on $[\mathcal{K}]$ rather than individual controllers.

Fixed-point detection (\texttt{python/check\_fixed\_point.py}) identifies when a
controller transformation leaves the controller unchanged within tolerance. Stability
accumulation tracks how consistently each class achieves fixed-point status. The
canonical representative of each class is the controller with highest stability score.

\subsection{The Fixed Point of the Total Operator}

The full system defines a composite restriction operator:
\[
\mathcal{C}_{\mathrm{total}} = \mathcal{R}_{13} \circ \cdots \circ \mathcal{R}_1
\]
Closure is achieved when:
\[
\mathcal{C}_{\mathrm{total}}(X^*) = X^*
\]
At this point all cocycles in $H^1(\mathcal{U}, \mathcal{S})$ are trivial, all type
mismatches are resolved by admissible bridge morphisms, all operator gaps are filled by
generated or retrieved morphisms, and all controller redundancies are collapsed in the
quotient. In terms of Section~\ref{sec:feasible}:
\[
\dim_{\mathrm{eff}}(\mathcal{M}_{13}) = 0
\]
The system is not a builder, not a learner, not a simulator. It is a
constraint-saturating operator whose sole function is to eliminate all inconsistent
configurations.

\subsection{Correspondence to the Main Theorems}

The filtration $\mathcal{M}_0 \supset \cdots \supset \mathcal{M}_{13}$ is the
implementation-level instantiation of the feasible set contraction established in
Sections~\ref{sec:identifiability} and~\ref{sec:disambiguation}. The main paper proves
that the consistency operator $\mathcal{C}$ has a unique fixed point when the projection
family is injective on $\mathcal{F}$ and dynamics separate residual ambiguity. The
constellation construction proves the same result one filtration level at a time:

Stages 1--4 correspond to the static identifiability condition ($\Omega_{\mathrm{obs}}$
contraction by sensor coverage). Stages 5--10 correspond to the dynamical disambiguation
condition ($\Gamma$ refinement by trajectory constraints). Stages 11--12 correspond to
the type-theoretic and transfer conditions that ensure the operator algebra is internally
consistent and portable. Stage 13 corresponds to the controller-level analogue of closure
in $\Xmod$: the decision procedure itself is fixed, not just the field state.

The fixed point $\mathcal{C}_{\mathrm{total}}(X^*) = X^*$ is therefore the operational
realization of the Yarncrawler Identifiability Theorem in a system of hundreds of
interconnected repositories. The constellation achieves closure when no alternative module
configuration, operator sequence, or controller policy remains consistent with all imposed
constraints simultaneously. Stages 14--16, which follow, extend the filtration from logical
closure to viability closure, from internal closure to external legibility, and from task
closure to mission closure.

\subsection{Stage 14: Homeostatic Constraint (Internal Sensory Monitoring)}

De Gyurky's Sensory System is responsible not only for external environmental awareness
but for constant monitoring of the system's internal health. In the filtration, this
stage introduces a viability requirement: the system may only occupy configurations that
do not threaten its own capacity for continued operation.

Let $\mathcal{S}_{\mathrm{infra}}$ denote the state space of the constellation's
infrastructure --- memory consumption, build-time latency, dependency graph depth,
repository count, and inter-service coupling density. Define a health functional
\[
\mathcal{H} : \mathcal{M}_{13} \to \mathbb{R}
\]
measuring departure from operational viability. The filtration step is:
\[
\mathcal{M}_{14} = \{ X \in \mathcal{M}_{13} \mid \mathcal{H}(X) \le \mathcal{H}_{\max} \}
\]
where $\mathcal{H}_{\max}$ is the viability threshold. This constraint prunes from the
feasible set any configuration that satisfies logical and type consistency but exceeds
resource bounds: a proposed repair that creates a dependency cycle exceeding memory
limits, or a composition chain that degrades build latency below operational thresholds,
is rejected even if it would satisfy all earlier filtration conditions.

The effect of Stage 14 is to make the Will stratum subordinate to the Sensory stratum.
Deployment is not only gated on logical closure; it is gated on metabolic stability.
A constellation that achieves $H^1 = 0$ but collapses under its own resource load has
not achieved closure in the full sense. $\mathcal{M}_{14}$ is the submanifold of
$\mathcal{M}_{13}$ on which the system can sustain itself.

In the language of the paper's main framework, Stage 14 adds an additional component to
the regularizer:
\[
\mathcal{R}_{14}(X) = \mathcal{R}(X) + \lambda_{\mathrm{health}}\,\mathcal{H}(X)
\]
so that configurations near the viability boundary are penalized even when logically
admissible. The admissible set $\mathcal{A}$ is refined to $\mathcal{A}_{14} = \mathcal{A}
\cap \mathcal{M}_{14}$.

\subsection{Stage 15: Presentational Projection (Observer-Accessible Interface)}

De Gyurky's Presentation System governs the transformation of internal experience into
projectable form: viewing, rendering, and making legible the system's internal state to
external observers. In the filtration, this stage introduces the requirement that the
system's closed internal state must admit a coherent external representation.

Define the interface space $\mathcal{I}$ as the space of observer-accessible artifacts:
documentation, visualizations, API specifications, structured logs, and public repository
metadata. Define a projection operator:
\[
\mathcal{P} : \mathcal{A}_{14} \to \mathcal{I}
\]
The filtration step is:
\[
\mathcal{M}_{15} = \{ X \in \mathcal{M}_{14} \mid \mathcal{P}(X) \approx y_{\mathrm{obs}} \}
\]
where $y_{\mathrm{obs}}$ is the target observation --- the mission requirement expressed
in the interface space. This is the presentational analogue of the sensor projection
condition $\|\Pi_i(X) - y_i\| \le \epsilon_i$: the system must not only be internally
consistent but must project into the interface space in a way that satisfies the
requirements of its external observers.

This stage formalizes the distinction between \emph{internal completion} (the collapsed
manifold $\mathcal{M}_{14}$) and \emph{external completion} (the legible projection
$\mathcal{P}(X^*)$). A constellation that has achieved internal closure but cannot
explain its own state --- cannot produce documentation from its registry, cannot render
its dependency graph, cannot expose its sensor projections as human-readable artifacts ---
has not achieved full closure. $\mathcal{M}_{15}$ is the submanifold of configurations
that are both internally solved and externally legible.

The Presentation System in the implementation corresponds to the \texttt{docs/generators/}
layer: scripts that generate FRD, FDD, SISD, and STIP documents automatically from the
registry, README stubs from module metadata, and interface specifications from type
signatures. These are not decorative outputs; they are constraints in $\mathcal{M}_{15}$
that must be satisfied before Will executes.

\subsection{Stage 16: The Thought System and Terminal Closure}

The Thought System --- the \emph{Nexus Cogitationis} --- is the hub of all communication
and the locus of self-awareness in de Gyurky's architecture. It is responsible for
detecting not merely that a specific task is closed, but that the entire constellation
has reached a state of structural closure across all strata simultaneously. In the
filtration, this is the terminal stage.

Define an equivalence relation $\sim_{\mathrm{mission}}$ on $\mathcal{M}_{15}$: two
configurations are mission-equivalent if they produce the same strategic outcome under
the constellation's governing objectives. Formally:
\[
X \sim_{\mathrm{mission}} X'
\;\Longleftrightarrow\;
\mathcal{P}(X) = \mathcal{P}(X') \;\text{ and }\;
\mathcal{H}(X) \approx \mathcal{H}(X')
\;\text{ and }\;
[\mathcal{K}(X)] = [\mathcal{K}(X')]
\]
where $[\mathcal{K}]$ denotes the controller equivalence class from Stage 13. The terminal
object is the quotient:
\[
\mathcal{M}_{16} = \mathcal{M}_{15} / \sim_{\mathrm{mission}}
\]

At this stage the system achieves what the main paper calls closure and what de Gyurky
calls genuine autonomy: the generator and the trajectory are indistinguishable, the
internal constraint structure and the external mission requirement are unified, and the
system operates at the fixed point of its own constraints rather than searching for it.
The Thought System acknowledges the global section's existence across all sixteen strata
and stabilizes the manifold.

The full filtration is:
\[
\mathcal{M}_0 \supset \mathcal{M}_1 \supset \cdots \supset \mathcal{M}_{13}
\supset \mathcal{M}_{14} \supset \mathcal{M}_{15} \supset \mathcal{M}_{16}
\]
and the composite restriction operator is:
\[
\mathcal{C}_{\mathrm{total}} = \mathcal{R}_{16} \circ \cdots \circ \mathcal{R}_1
\]
Closure is:
\[
\mathcal{C}_{\mathrm{total}}(X^*) = X^*, \qquad \dim_{\mathrm{eff}}(\mathcal{M}_{16}) = 0
\]

At the fixed point: all cohomological obstructions are trivialized ($H^1 = 0$), all type
mismatches are resolved, all operator gaps are filled, all controller redundancies are
collapsed, the system is metabolically viable ($\mathcal{H}(X^*) \le \mathcal{H}_{\max}$),
the internal state is externally legible ($\mathcal{P}(X^*) \approx y_{\mathrm{obs}}$),
and the mission equivalence class is a singleton ($|\mathcal{M}_{16}| = 1$).

This is the mathematical definition of a done autonomous state. The constellation is not
working toward a goal; it is the realized solution, operating at the fixed point of its
own internal and external constraints across all sixteen strata of the filtration.

%% ============================================================
\section{The Constellation Architecture as Operational Sheaf}
\label{app:constellation}
%% ============================================================

\subsection{From Deployment Tool to Sheaf Realization}

The implementation described in Appendix~\ref{app:implementation} is not a deployment
pipeline that happens to use sheaf language. It is a direct computational realization
of the sheaf $\mathcal{S}_\Pi$ over the quotient configuration space $\Xmod$. The
correspondence is exact, not analogical.

Each software module is a local section $s_i \in \mathcal{S}_\Pi(U_i)$: it has a
declared type signature (the \texttt{inputs} and \texttt{outputs} fields of
\texttt{module.yaml}), a declared stratum (one of the eight de Gyurky functional roles),
and an explicit dependency list (the \texttt{depends\_on} field). The dependency list
defines the restriction maps $\rho_{ij}$: module $j$ depends on module $i$ means $i$'s
output must be compatible with $j$'s input, i.e., $\rho_{ij}(s_i) = s_j|_{U_i \cap U_j}$.

The \texttt{check\_commutativity.sh} script detects violations of the gluing condition.
When module $j$ requires field $\Phi$ in space $H^1(\Omega)$ but module $i$ provides
$\Phi$ in space $L^2(\Omega)$, the script emits a structured object of the form
\begin{verbatim}
{
  "type": "cohomology_obstruction",
  "from": "module_i", "to": "module_j",
  "field": "Phi"
}
\end{verbatim}
This is not an error message in the conventional sense. It is a representation of a
non-trivial \v{C}ech 1-cocycle $c_{ij} \in H^1(\mathcal{U}, \mathcal{S}_\Pi)$. The
system has computed that the local sections $s_i$ and $s_j$ do not agree on the
overlap $U_i \cap U_j$.

Successful deployment --- the condition under which \texttt{run.sh} completes without
error and all repositories are materialized via \texttt{gh repo create} --- corresponds
exactly to the existence of a unique global section: $H^1(\mathcal{U}, \mathcal{S}_\Pi)
= 0$ and $\dim H^0 = 1$. GitHub becomes the observable projection of a globally
consistent semantic object, not the ground truth.

\subsection{The Eight Strata as Factorization of $\mathcal{C}$}

The eight de Gyurky systems --- Sensory, Intellect, Understanding, Decision, Reason,
Will, Presentation, Thought --- are not implemented as separate components. They are
implemented as functional strata through which every module is classified, and together
they constitute a factorization of the consistency operator $\mathcal{C}$.

The \emph{Sensory} stratum defines the projection operators $\{\Pi_i\}$: it comprises
the modules that produce observations from field configurations. Every module in this
stratum has typed outputs corresponding to sensor channels (infrared, LiDAR, RGB, audio,
photonic strip) and its \texttt{module.yaml} declares the field space from which it
projects.

The \emph{Intellect} stratum corresponds to the encoder $q_\phi$: it comprises the
modules that generate candidate field configurations from observations, operating in the
high-entropy regime without enforcing admissibility. This is the proposal mechanism of
Section~\ref{sec:inversion}.

The \emph{Understanding} stratum normalizes and packages candidate configurations for
the consistency loop. In sheaf terms, it constructs the overlap data: given two locally
consistent configurations from different sensors, it prepares the restriction maps for
comparison.

The \emph{Reason} stratum is the admissibility filter $\mathcal{R}$: modules in this
stratum enforce that field configurations satisfy the RSVP equations, ecological
plausibility constraints, and civilizational feasibility. A configuration that fails
Reason is not merely penalized --- it is excluded from $\mathcal{A}$ entirely.

The \emph{Decision} stratum implements operator selection. At each iteration of the
refinement loop, Decision modules select among candidate morphisms by evaluating them
against the current state and observations. This is the variational minimization in
Definition~5.3, made concrete as a selection procedure over a ranked candidate pool.

The \emph{Will} stratum executes the selected morphism: it applies the chosen operator
to the current field configuration and, when closure is certified, materializes the
result into the external world (repository creation, push, tagging). Critically, Will
modules are never invoked until the closure check passes. Deployment is not a routine
step in the pipeline; it is the consequence of $|\mathcal{F}| = 1$.

The \emph{Presentation} stratum projects the internal field state back into observable
form: documentation generation, README synthesis, visual output. It corresponds to the
forward model in the refinement loop (Section~\ref{sec:inversion}), applied not for
validation but for communication.

The \emph{Thought} stratum is the consistency operator itself: the orchestration loop
in \texttt{run.sh} that iterates over all strata in the correct order, checks closure,
triggers repair when it fails, and halts when the global section is achieved. Thought
does not have its own modules; it is the loop that runs them.

The complete factorization is:
\[
\mathcal{C} = \mathrm{Will} \circ \mathrm{Decision} \circ \mathrm{Understanding}
\circ \mathrm{Intellect} \circ \mathrm{Sensory} \circ \mathrm{Reason}
\]
with Presentation as an output projection and Thought as the iteration of this
composition until fixed-point.

\subsection{The Repair Loop as Coboundary Construction}

The three-script repair loop --- \texttt{propose\_bridge.sh},
\texttt{evaluate\_bridge.sh}, \texttt{apply\_repair.sh} --- implements the sheaf repair
operator $\mathcal{R}_{\mathrm{sheaf}}$ of Section~\ref{sec:equivalence} in executable
form.

When \texttt{check\_commutativity.sh} detects a cohomological obstruction between
modules $i$ and $j$ over field $\Phi$ in incompatible spaces, \texttt{propose\_bridge.sh}
reads the structured obstruction object and queries the type graph defined in
\texttt{schema/typing\_rules.yaml}. The type graph is a directed acyclic graph over
field spaces: $H^1(\Omega) \hookrightarrow L^2(\Omega) \hookrightarrow \mathrm{Scalar}$.
The path resolver (\texttt{engine/type\_path\_resolver.sh}) performs a breadth-first
search over this graph to find an admissible sequence of embeddings from the provided
space to the required space.

Each step of the resolved path generates one bridge module. If the path is
$H^1(\Omega) \to L^2(\Omega) \to \mathrm{Scalar}$, then on the first repair pass the
system generates and inserts:
\[
\texttt{phi\_bridge\_H1\_to\_L2} : \Phi_{H^1} \mapsto \Phi_{L^2}
\]
On the second pass, module $j$ still lacks the Scalar-type $\Phi$ it requires, but
now a provider exists in $L^2(\Omega)$, and the system generates:
\[
\texttt{phi\_bridge\_L2\_to\_Scalar} : \Phi_{L^2} \mapsto \Phi_{\mathrm{Scalar}}
\]
This incremental coboundary construction is precisely the multi-hop path completion
established in Section~\ref{sec:equivalence}: each bridge module trivializes one edge
of the \v{C}ech cocycle, and the loop continues until $H^1 = 0$.

The loop is bounded by \texttt{MAX\_ITERS} in \texttt{run.sh}. If closure is not
achieved within that bound, the system emits a structured failure report identifying
which cocycles remain non-trivial and whether the failure is projection degeneracy
(the type graph contains no path between the required spaces) or regularizer flatness
(multiple valid bridge paths exist and the system cannot select among them without
model revision).

%% ============================================================
\section{The Morphism Kernel as RSVP Operator Algebra}
\label{app:morphism_kernel}
%% ============================================================

\subsection{Operators as RSVP Field Transformations}

The morphism kernel implements the operator class library of Appendix~\ref{app:operators}
as an executable algebra over the RSVP field triple. Every morphism in the system ---
whether symbolic, learned, hybrid, composed, or generated --- is normalized into a
common affine RSVP form:
\[
(\Phi, \mathbf{v}, S) \;\mapsto\;
(a_\Phi\, \Phi + b_\Phi,\;\; a_v\, \mathbf{v} + \mathbf{b}_v,\;\; a_S\, S + b_S)
\]
where the six parameters $(a_\Phi, b_\Phi, a_v, \mathbf{b}_v, a_S, b_S)$ constitute the
operator's normal form. This normalization, implemented in \texttt{python/morphism\_normalize.py},
is what makes composition exact rather than approximate: the normal form of a composition
$M_2 \circ M_1$ is computed analytically by substitution, so no numerical error
accumulates through chains of operators.

The normalization also makes blending in operator space well-defined. A context-weighted
blend of operators $\{M_k\}$ with weights $\{w_k\}$ is the operator with parameters:
\[
a_\Phi^{\mathrm{blend}} = \sum_k w_k a_\Phi^{(k)}, \quad
b_\Phi^{\mathrm{blend}} = \sum_k w_k b_\Phi^{(k)}, \quad \ldots
\]
This is implemented in \texttt{python/blend\_morphisms.py} and used both for
context-neighbor retrieval (Section~\ref{app:morphism_kernel}.3 below) and for
multi-controller blending (Appendix~\ref{app:controller_quotient}).

The RSVP field triple enters the implementation not as three independent channels but as
a coupled physical object. The vector field $\mathbf{v}$ is implemented as edge-based
flux on an explicit graph topology (\texttt{state/topology.json}): each edge carries a
scalar flux value, and divergence at node $i$ is computed as the exact discrete
conservation law $\mathrm{div}(i) = \sum_{e \to i} f_e - \sum_{i \to e} f_e$. The
entropy field $S$ is updated by the discrete continuity equation accounting for local
divergence and shear stress along edges. This means the three fields are genuinely
coupled in the forward dynamics, not merely co-present.

\subsection{Sensor-Coupled Operator Selection}

Morphism selection is not a combinatorial search over operator space. It is a
variational minimization over a candidate pool, evaluated against sensor observations
and temporal consistency requirements. The scoring functional is:
\[
\mathcal{J}(M) = w_1 \mathcal{L}_{\mathrm{now}}(M)
+ w_2 \mathcal{L}_{\mathrm{future}}(M)
+ w_3 \mathcal{P}_{\mathrm{div}}(M)
+ w_4 \mathcal{P}_{\mathrm{boundary}}(M)
+ w_5 \mathcal{P}_{\mathrm{rough}}(M)
- \mathcal{B}_{\mathrm{semantic}}(M)
+ \mathcal{P}_{\mathrm{semantic}}(M)
\]
where $\mathcal{L}_{\mathrm{now}}$ is the projection residual at the current step,
$\mathcal{L}_{\mathrm{future}}$ is the residual after one forward simulation step,
$\mathcal{P}_{\mathrm{div}}$ penalizes unphysical divergence, $\mathcal{P}_{\mathrm{boundary}}$
penalizes flux leakage through boundary nodes, $\mathcal{P}_{\mathrm{rough}}$ penalizes
rapid spatial variation in edge flux, and $\mathcal{B}_{\mathrm{semantic}}$ and
$\mathcal{P}_{\mathrm{semantic}}$ are bonuses and penalties based on whether the operator's
declared semantic properties (e.g., \texttt{divergence\_preserving},
\texttt{entropy\_relaxing}, \texttt{conservative}) are actually realized in the field
transformation it produces.

The semantic property system is the key upgrade over standard variational selection. An
operator that declares itself \texttt{entropy\_relaxing} but produces increasing mean
entropy is penalized; one that achieves entropy reduction when it claimed none gets no
bonus. This couples the operator's declared intention to its empirical behavior, making
the selection criterion self-validating. In the language of the paper's main theorems,
this enforces that the regularizer $\mathcal{R}$ is evaluated not only as an
admissibility condition on field configurations but as a consistency condition on
operator semantics.

Scoring is extended to trajectory rollouts: rather than evaluating at a single step, the
system simulates a morphism forward over a horizon $T$ (typically 8--10 steps) and
integrates the scoring functional over the resulting trajectory. This implements
counterfactual consistency (Section~\ref{sec:closure}) as the primary criterion for
morphism admission: an operator is selected not because it fits the current observation
but because it induces a trajectory that remains within tolerance of observations across
the full horizon.

\subsection{Context-Sensitive Memory and Operator Discovery}

The morphism registry stores operators not as a flat table but as a context-indexed
statistical manifold. Each entry carries:
\begin{itemize}
\item the normalized RSVP operator parameters,
\item declared semantic properties,
\item usage count and mean/variance of rollout scores,
\item a context centroid computed as the average context vector across all uses.
\end{itemize}
Retrieval is context-weighted: given the current field configuration and sensor
observations, a context vector is computed (mean $\Phi$, mean flux, mean entropy, sensor
targets), and candidate operators are ranked by a combined score:
\[
\mathrm{rank}(M) = \bar{s}(M) - \lambda\, d(\mathbf{c}, \bar{\mathbf{c}}(M))
\]
where $\bar{s}$ is mean rollout score, $\mathbf{c}$ is the current context, and
$\bar{\mathbf{c}}(M)$ is the operator's context centroid. Operators are therefore
preferred both when they have historically performed well and when the current context
resembles the contexts in which they were used. This implements a minimal version of the
contextual retrieval that distinguishes Yarncrawler from global optimization.

When no operator in the registry achieves sufficient score and no admissible type path
exists between the required field spaces, the system invokes morphism synthesis: it
generates symbolic candidate operators of declared types (projection, aggregation,
embedding) and evaluates them against the sensor-coupled scoring functional. Candidates
that pass the admission threshold (\texttt{python/admit\_morphism.py}) enter the registry
and become available for future retrieval and composition.

Composition discovery operates over the history log \texttt{logs/history.jsonl}: pairs
of operators that were applied in sequence with consistently high rollout scores are
identified as candidate compositions (\texttt{python/discover\_compositions.py}) and
stored as composite morphisms. When a composite achieves high score across many distinct
contexts, it is abstracted into a single atomic operator (\texttt{python/abstract\_morphism.py}),
compressing the composition chain into a reusable primitive. This is the operational
version of the operator abstraction described in the main text: just as the operator
class library grows by extending $\Theta$ with new process generators $F_k$, the
implementation grows its registry by abstracting empirically successful operator chains.

Cross-domain transfer (\texttt{python/transfer\_morphism.py}) applies the functor
$\Psi$ of Theorem~\ref{thm:equivalence} in the opposite direction: given a source domain
descriptor, a target domain descriptor, and a transfer map, it transports a registry
operator from one sensor bundle and topology into another. The transport scales
bias parameters by the ratio of node and edge counts, maps field type signatures through
the declared type map, and annotates the resulting operator with a functoriality score
computed by \texttt{python/check\_functoriality.py}. Transferred operators are admitted
to the target registry with a score discount proportional to $(1 - s_F)$, reflecting
the theoretical result that transfer is approximate rather than exact unless the transfer
map is a full isomorphism.

%% ============================================================
\section{The Controller Quotient and Meta-Level Closure}
\label{app:controller_quotient}
%% ============================================================

\subsection{Controllers as Functors}

Each controller --- beam search, dynamic programming, learned gating, spatial gating,
repair --- is a functor from the category of observable states to the category of
candidate morphism sequences:
\[
\mathcal{K} : (\text{state}, \text{registry}) \;\longrightarrow\; \text{candidate sequence}
\]
The five controllers in \texttt{controllers/registry/} correspond to five distinct
functorial strategies for generating morphism sequences given the current field state and
observation. Each strategy reflects a different epistemic regime: beam search when
uncertainty is high and exploration is warranted, dynamic programming when the regime is
stable enough to admit policy decomposition, learned gating when the current context is
well-represented in the registry and amortized inference is reliable, spatial gating when
different regions of the field require different operator policies, and repair when an
explicit obstruction has been detected and a targeted coboundary must be constructed.

The five controllers do not simply select among these strategies by hard switching.
They are blended continuously. The controller confidence layer
(\texttt{python/controller\_confidence.py}) computes, at each time step, a weight for
each controller based on: the margin between best and second-best candidates (beam
confidence), the regime stability score (DP confidence), and the context-similarity of
the current state to well-populated registry regions (learned confidence). These weights
are fed into \texttt{python/blend\_controllers.py}, which normalizes the controllers'
selected operators into a common RSVP affine form and computes a weighted blend. The
resulting operator inherits the property union of its contributing controllers and is
applied as a single morphism at that time step.

At a finer spatial scale, the spatial controller layer extends this blending to be
node-specific. Controller weights are computed per node based on local field properties
(entropy, flux magnitude, boundary status) and smoothed over the graph topology via a
discrete Laplacian. The resulting node-wise weight field
(\texttt{python/spatially\_smooth\_weights.py}) is applied in
\texttt{python/blend\_nodewise\_controllers.py} to produce an operator that varies
spatially: boundary nodes may receive high DP weight (conservative, boundary-respecting),
high-entropy interior nodes may receive high beam weight (exploratory), and
low-entropy familiar regions may receive high learned weight (amortized). The spatial
controller is therefore not a single morphism but a field of morphisms over the topology,
implemented as the nodewise hybrid type in \texttt{python/morphism\_apply.py}.

\subsection{Natural Transformations as Controller Evolution}

A natural transformation between controllers is a mapping
$\eta : \mathcal{K}_1 \Rightarrow \mathcal{K}_2$ that translates one controller's
parameters and decision structure into another's. The implementation stores these as
JSON objects in \texttt{controllers/transformations/}: for example,
\texttt{beam\_to\_learned.json} maps beam search parameters (width, depth) to learned
controller parameters (attention span, planning horizon) with an explicit weight
controlling how strongly the transformation is applied.

Controller transformation is applied iteratively. At each step, the current controller
state $\mathcal{K}_t$ is fed through the transformation to produce $\mathcal{K}_{t+1}$
(\texttt{python/iterate\_controller.py}), and the distance between $\mathcal{K}_t$ and
$\mathcal{K}_{t+1}$ is computed (\texttt{python/controller\_distance.py}). When the
distance falls below a threshold $\varepsilon$, a fixed-point signal is emitted
(\texttt{python/check\_fixed\_point.py}). The stability score of the controller is
incremented when a fixed point is detected and decayed otherwise
(\texttt{python/update\_stability.py}). This implements, at the controller level, the
same fixed-point iteration that the main consistency operator $\mathcal{C}$ performs at
the field level.

The formal analogy is precise. The field-level fixed point is:
\[
\mathcal{C}([X^*]) = [X^*]
\]
The controller-level fixed point is:
\[
\eta(\mathcal{K}^*) \approx \mathcal{K}^*
\]
Both are instances of the same underlying principle: a system has closed when applying
its own transformation leaves it unchanged. The field-level closure corresponds to
$|\mathcal{F}| = 1$ and the theorems of Sections~\ref{sec:identifiability}
and~\ref{sec:disambiguation}. The controller-level closure corresponds to a stable
decision procedure that no transformation meaningfully alters.

Transformations can be learned from the controller trace log
(\texttt{logs/controller\_trace.jsonl}) by identifying which controller decisions, in
which contexts, led to high rollout scores (\texttt{python/learn\_controller\_transform.py}).
This makes controller evolution data-driven rather than hand-specified: the system
discovers which transformations improve decision quality and encodes them as natural
transformations in the registry.

\subsection{The Quotient Structure over Controller Space}

Two controllers are observationally equivalent if their decision behavior is
indistinguishable across all states in the registry: they select the same morphisms
(or morphisms with equivalent rollout scores) in all contexts where both have been
applied. This is measured by the controller distance metric, and the equivalence
relation is the thresholded version: $\mathcal{K}_1 \sim \mathcal{K}_2$ iff
$d(\mathcal{K}_1, \mathcal{K}_2) < \varepsilon$.

The quotient $\mathrm{Controllers}/{\sim}$ is computed by
\texttt{python/build\_controller\_quotient.py}: it clusters all controllers into
equivalence classes and selects a canonical representative per class
(\texttt{python/select\_representatives.py}, choosing the controller with highest
stability score within each class). The resulting structure is stored in
\texttt{controllers/quotient/}, and all subsequent controller selection operates over
this quotient rather than over raw controllers.

This quotient structure is the meta-level counterpart of the quotient $\Xmod$ that
governs field-state identifiability. Just as two field configurations that agree on all
sensor projections are identified in $\Xmod$ --- because the system cannot and should not
distinguish them --- two controllers that agree on all decision outputs across all
observable contexts are identified in $\mathrm{Controllers}/{\sim}$. The system operates
on equivalence classes of decision procedures, not on individual procedures.

The consequence is a compression of the decision space that mirrors the compression of
the field space. Just as the feasible set $\mathcal{F}$ contracts toward a singleton
under increasing sensor coverage and trajectory constraints, the quotient controller
space contracts toward a minimal set of genuinely distinct decision structures as the
system accumulates experience. Controllers that appeared distinct during early exploration
may prove equivalent once the registry is populated, because their differences are only
visible in contexts that the system has not yet encountered.

\subsection{Closure at the Meta-Level}

The system achieves meta-level closure when:
\begin{enumerate}
\item The morphism registry has stabilized: new operators are admitted rarely, and the
candidate pool for any new context is dominated by retrieved operators rather than
freshly generated ones.
\item The controller quotient has stabilized: the equivalence classes no longer split
or merge as new controller traces are added, and the set of canonical representatives
is fixed.
\item The transformation graph between controllers is dominated by intra-class
transformations: most observed controller evolutions map a controller to an equivalent
one, with inter-class transitions becoming rare.
\end{enumerate}
These three conditions are the meta-level analogue of the three closure criteria of
Section~\ref{sec:closure}: projection parity (sensor residuals within tolerance),
dynamical stability (field state on an admissible trajectory), and counterfactual
consistency (forward simulation agrees with future observations). At meta-level closure,
the system has not merely found a world-state consistent with all observations --- it has
found a decision procedure consistent with all operating contexts. The architecture is
not just closed with respect to the world; it is closed with respect to itself.

This self-referential closure is the operational realization of a principle that runs
throughout the framework: a system is done not when it has found a solution but when
applying its own processes to that solution leaves it unchanged. In the language of
sheaves, the global section exists not only over the cover of sensor modalities but over
the cover of decision contexts. The Yarncrawler constellation does not merely reconstruct
a world-state from observations; it identifies the invariant decision structure that
makes reconstruction reliable across all the worlds it might encounter.

\newpage
\bibliographystyle{plainnat}
\begin{thebibliography}{99}

\bibitem{friston2010}
Friston, K. (2010).
The free-energy principle: a unified brain theory?
\textit{Nature Reviews Neuroscience}, 11(2), 127--138.

\bibitem{friston2019}
Friston, K., FitzGerald, T., Rigoli, F., Schwartenbeck, P., \& Pezzulo, G. (2019).
Active inference: a process theory.
\textit{Neural Computation}, 31(1), 1--49.

\bibitem{hordijk2019}
Hordijk, W., \& Steel, M. (2019).
Autocatalytic networks: an intimate relation between network topology and dynamics.
\textit{Bulletin of Mathematical Biology}, 81(5), 1357--1374.

\bibitem{maturana1980}
Maturana, H. R., \& Varela, F. J. (1980).
\textit{Autopoiesis and Cognition: The Realization of the Living}.
D.~Reidel.

\bibitem{rosen1985}
Rosen, R. (1985).
\textit{Anticipatory Systems}.
Pergamon Press.

\bibitem{abramsky2011}
Abramsky, S., \& Brandenburger, A. (2011).
The sheaf-theoretic structure of non-locality and contextuality.
\textit{New Journal of Physics}, 13, 113036.

\bibitem{spivak2014}
Spivak, D. I. (2014).
\textit{Category Theory for the Sciences}.
MIT Press.

\bibitem{lewis2020}
Lewis, P., Perez, E., et al. (2020).
Retrieval-augmented generation for knowledge-intensive NLP tasks.
\textit{Advances in Neural Information Processing Systems}, 33.

\bibitem{mildenhall2021}
Mildenhall, B., Srinivasan, P. P., Tancik, M., Barron, J. T., Ramamoorthi, R.,
\& Ng, R. (2021).
NeRF: representing scenes as neural radiance fields for view synthesis.
\textit{Communications of the ACM}, 65(1), 99--106.

\bibitem{freidlin1984}
Freidlin, M. I., \& Wentzell, A. D. (1984).
\textit{Random Perturbations of Dynamical Systems}.
Springer.

\end{thebibliography}

\end{document}
